\documentclass[11pt]{amsart}

\usepackage[a4paper,margin=2.5cm]{geometry}
\usepackage{amsmath,amssymb,amsthm,mathtools}
\usepackage{enumitem}
\usepackage{setspace}
\usepackage[hidelinks]{hyperref}

\setstretch{1.16}

\newtheorem{theorem}{Theorem}[section]
\newtheorem{proposition}[theorem]{Proposition}
\newtheorem{lemma}[theorem]{Lemma}
\newtheorem{corollary}[theorem]{Corollary}
\newtheorem{claim}[theorem]{Claim}
\theoremstyle{definition}
\newtheorem{definition}[theorem]{Definition}
\newtheorem{problem}[theorem]{Problem}
\newtheorem{remark}[theorem]{Remark}

\newcommand{\R}{\mathbb R}
\newcommand{\HH}{\mathcal H}
\newcommand{\Hn}{\mathcal H^{n-1}}
\newcommand{\eps}{\varepsilon}
\newcommand{\Om}{\Omega}
\newcommand{\Per}{P}
\newcommand{\dx}{\,dx}
\newcommand{\dH}{\,d\mathcal H^{n-1}}
\newcommand{\ds}{\,d\sigma}
\newcommand{\abs}[1]{\left|#1\right|}
\newcommand{\norm}[1]{\left\|#1\right\|}
\newcommand{\osc}{\mathcal O}
\newcommand{\Def}{\mathcal D}
\newcommand{\Dev}{\mathcal E}
\DeclareMathOperator{\barint}{\mathop{\ooalign{$\int$\cr\hidewidth--\hidewidth\cr}}}

\title[Attempts at an effective Fusco--Julin inequality]
{Attempts at an effective Fusco--Julin\\
strong quantitative isoperimetric inequality}
\author{}
\date{May 2026}

\begin{document}

\begin{abstract}
The strong quantitative isoperimetric inequality of Fusco and Julin controls
a same-centre sum of volume asymmetry and radial normal oscillation by the
perimeter deficit; its proof reduces this to a bound for the normal
oscillation index \(\beta\).  The published proof gives a dimensional constant, but not a
computable formula for that constant.  This article collects rigorous attempts
to make the argument effective in the sense of certifying computability of
all dimensional constants, without evaluating them numerically.  First, we establish all normalisations and
the radial Gauss--Green identity in the form needed for later use.  Second, we
give a complete explicit theorem for balanced nearly-spherical sets: if a
boundary is a sufficiently small \(W^{1,\infty}\) radial graph over a sphere,
then one fixed centre controls both its symmetric difference and its radial
normal oscillation by the perimeter deficit, with displayed constants.  The
quadratic constant for the oscillation term is \(2n/(n+1)\), dictated by the
degree-two spherical harmonics.  Third, we first test a Plan~1-style
decomposition \(\beta^2=\mathcal D+\mathcal C\), then implement the
regularising Fusco--Julin selection which penalises \(\beta^2\) itself.  For
the selected minimiser we prove, with displayed constants, coarse area
quasi-minimality, density, localisation of all radial-potential maximisers,
and additive perimeter almost-minimality after annular trapping.  Finally, we
give a computability certificate for the standard graph-entry step: the
compactness-based harmonic approximation in the Tamanini--Maggi proof is
replaced by a finite spectral truncation argument.  This closes the
regularised selection route on the bounded class with computable
constants.  Finally, the finite slicing reduction already used by Fusco and
Julin is certified computable and composed with that bounded closure.  Their
constructive comparison between the normal index and the same-centre combined
index then yields the original Fusco--Julin theorem for arbitrary
finite-perimeter sets with a computable dimensional constant.
\end{abstract}

\maketitle
\tableofcontents

\section{Aim, status, and conventions}\label{sec:intro}

\subsection{The theorem whose constant is to be made effective}

Throughout the article \(n\ge2\), \(B_\rho(z)\) denotes the open Euclidean ball
of radius \(\rho\) and centre \(z\), \(B_\rho=B_\rho(0)\), and
\(\omega_n=|B_1|\).  If \(E\subset\R^n\) has finite perimeter, we write
\(\partial^*E\) for its reduced boundary and \(\nu_E\) for its
measure-theoretic outer unit normal.

\begin{definition}[Radial normal oscillation]\label{def:beta}
Let \(E\subset\R^n\) be a set of finite perimeter with
\(|E|=|B_\rho|\).  For \(z\in\R^n\), set
\[
   e_z(x):=\frac{x-z}{|x-z|}
   \quad\text{for }x\ne z,
\]
and
\[
   \osc_z(E):=
   \frac12\int_{\partial^*E}|\nu_E-e_z|^2\,d\Hn.
\]
The Fusco--Julin oscillation index is
\[
   \beta(E)^2:=\inf_{z\in\R^n}\osc_z(E).
\]
The unnormalised perimeter deficit is
\[
   \Def(E):=\Per(E)-\Per(B_\rho).
\]
\end{definition}

The reduced normal-oscillation target in the proof of Fusco and Julin
\cite{FJ} is the following.

\begin{theorem}[Normal-oscillation form of Fusco--Julin]\label{thm:FJ}
For every \(n\ge2\), there exists \(C_{\rm FJ}(n)>0\) such that every
finite-perimeter set \(E\subset\R^n\) satisfying \(|E|=|B_\rho|\) obeys
\[
   \beta(E)^2\le C_{\rm FJ}(n)\Def(E).
\tag{1.1}\label{eq:FJ}
\]
\end{theorem}

Their published main theorem is stated for the same-centre combined
asymmetry \(\mathcal A_{\rm FJ}\) defined in
\eqref{eq:A-FJ-combined} below.  Proposition~1.2 of \cite{FJ} reduces
that statement to \eqref{eq:FJ}; Corollary~\ref{cor:original-FJ-computable}
records the effective recovery of their formulation.
The purpose here is narrower and more exact:
we seek a proof in which \(C_{\rm FJ}(n)\), and every smallness threshold used
to reach it, is certified computable from a finite chain of inequalities.  We
do not require the resulting numbers to be numerically evaluated.  The regularity-based
proof in \cite{FJ} proves existence of a dimensional constant, but does not
provide such a formula; see also the presentation of its argument in
\cite[Section 5.3]{FuscoSurvey}.

\subsection{What this article proves and what it does not prove}

There are four complete new pieces in this article.
Theorem~\ref{thm:graph-main} proves an explicit version of
\eqref{eq:FJ} for balanced \(W^{1,\infty}\)-small radial graphs.  The theorem
also controls the volume error and normal oscillation about the \emph{same}
centre, which is the form needed in the level-set moving-ball application.
Theorem~\ref{thm:FJ-selection} gives a Plan~1-style diagnostic selection
which preserves the quotient between perimeter deficit and radial potential
deficit; Proposition~\ref{prop:near-min} explains why it is not directly a
regularising selection.  Theorems~\ref{thm:true-beta-selection} and
\ref{thm:additive-almost-min} then implement the correct Fusco--Julin
regularisation: a minimiser preserving \(\Def/\beta^2\) is upgraded to an
additive perimeter almost-minimiser by quantitative annular trapping and
localisation of the potential centres.
The computability certificate in
Proposition~\ref{prop:flat-number-computable} then supplies the graph-entry
constant required for the bounded regularised-selection closure,
Proposition~\ref{prop:true-beta-bounded-close}.
Lemma~\ref{lem:computable-confinement} certifies computability of the
Fusco--Julin confinement step, and
Theorem~\ref{thm:global-computable-FJ} proves \eqref{eq:FJ} with a
computable dimensional constant.  Corollary~\ref{cor:original-FJ-computable}
then gives the original combined formulation with a computable constant.

There are two separate intrinsic attempts not closed here:
\begin{enumerate}[label=\textnormal{(\roman*)},leftmargin=2.2em]
\item Section~\ref{sec:reilly} derives a torsion/Reilly route relevant to
torsion superlevel sets.  The identities are proved, but their use for general
level sets requires a gradient-tail estimate and a bulk-to-boundary trace
estimate which are not proved here.
\item Section~\ref{sec:intrinsic} derives slicing and meridian-action
identities modelled on the geometric FMP proof and on coherent-level
arguments in functional Brunn--Minkowski stability.  It also proves the
centred-reflection reduction to \(n-1\) transverse symmetries and closes
the full oriented symmetry reduction in dimension \(2\).  For
\(n\ge3\), full centred symmetry, slice-motion control, and centred
meridian closure remain proposed problems, not proved inputs.
\end{enumerate}

\begin{remark}[Two corrections to preliminary notes]\label{rem:corrections}
Two numerical points must be fixed at the outset.  The Laplace--Beltrami
eigenvalue of degree-two spherical harmonics on \(S^{n-1}\) is
\[
   2(2+n-2)=2n,
\]
not \(2n+2\).  Consequently, under Definition~\ref{def:beta}, the quadratic
constant for degree-two balanced radial perturbations is
\[
   \frac{2n}{n+1};
\]
this is proved in Proposition~\ref{prop:degree-two}.  Also,
\texttt{arXiv:2410.20844} is not a Fusco survey; the source used below for a
survey account of the FJ proof is Fusco's 2015 article
\cite{FuscoSurvey}.
\end{remark}

\section{The radial identity and elementary reductions}\label{sec:identity}

This section spells out the identities which explain the choice of
\(\beta(E)\).  The calculation is elementary, but it is important that it be
written with the correct factor \(1/2\), the correct scaling, and no hidden
smoothness assumption.

\subsection{A finite-perimeter radial Gauss--Green identity}

\begin{lemma}[Radial divergence field]\label{lem:radial-div}
Fix \(z\in\R^n\), and define \(e_z(x)=(x-z)/|x-z|\) for \(x\ne z\).
Then
\[
   \operatorname{div}e_z(x)=\frac{n-1}{|x-z|}
   \qquad (x\ne z).
\tag{2.1}\label{eq:div-e}
\]
Moreover \(x\mapsto |x-z|^{-1}\) is locally integrable on \(\R^n\).
\end{lemma}

\begin{proof}
By translation it suffices to take \(z=0\).  For \(x\ne0\), write
\[
   e_0(x)=\frac{x}{|x|}.
\]
For each \(i\in\{1,\ldots,n\}\),
\[
   \partial_i\!\left(\frac{x_i}{|x|}\right)
   =\frac1{|x|}-\frac{x_i^2}{|x|^3}.
\]
Summing in \(i\) gives
\[
   \operatorname{div}e_0
   =\frac{n}{|x|}-\frac{|x|^2}{|x|^3}
   =\frac{n-1}{|x|}.
\]
Finally,
\[
   \int_{B_r(z)}\frac{dx}{|x-z|}
   =n\omega_n\int_0^r s^{n-2}\,ds
   =\frac{n\omega_n}{n-1}r^{n-1}<\infty,
\]
because \(n\ge2\).
\end{proof}

\begin{proposition}[Radial Gauss--Green formula]\label{prop:radial-GG}
Let \(E\subset\R^n\) be a bounded set of finite perimeter.  Then for every
\(z\in\R^n\),
\[
   \int_{\partial^*E}\nu_E\cdot e_z\,d\Hn
   =(n-1)\int_E\frac{dx}{|x-z|}.
\tag{2.2}\label{eq:radial-GG}
\]
Consequently,
\[
   \osc_z(E)
   =\Per(E)-(n-1)\int_E\frac{dx}{|x-z|}.
\tag{2.3}\label{eq:osc-potential}
\]
\end{proposition}

\begin{proof}
The field \(e_z\) is singular at \(z\), so we first remove that singularity.
Choose a smooth function \(\chi_\eps:[0,\infty)\to[0,1]\) satisfying
\[
   \chi_\eps(r)=0\quad (0\le r\le\eps),\qquad
   \chi_\eps(r)=1\quad (r\ge2\eps),\qquad
   |\chi_\eps'(r)|\le\frac{2}{\eps}.
\]
Set
\[
   X_\eps(x):=\chi_\eps(|x-z|)e_z(x).
\]
Since \(E\) is bounded, the finite-perimeter Gauss--Green formula, after an
irrelevant compactly supported cut-off outside a ball containing \(E\), gives
\[
   \int_{\partial^*E}X_\eps\cdot\nu_E\,d\Hn
   =\int_E\operatorname{div}X_\eps\,dx.
\tag{2.4}\label{eq:cutoff-GG}
\]
On \(\R^n\setminus\{z\}\),
\[
   \operatorname{div}X_\eps
   =\chi_\eps(|x-z|)\frac{n-1}{|x-z|}
     +\chi_\eps'(|x-z|).
\]
The second term is supported in
\(B_{2\eps}(z)\setminus B_\eps(z)\), and hence
\[
   \left|\int_E\chi_\eps'(|x-z|)\,dx\right|
   \le \frac2\eps |B_{2\eps}|
   =2^{n+1}\omega_n\eps^{n-1}\longrightarrow0.
\]
The first term converges in \(L^1(E)\) to
\((n-1)|x-z|^{-1}\) by Lemma~\ref{lem:radial-div} and dominated convergence.
On the boundary side, \(X_\eps\to e_z\) at every point other than \(z\), and
\(\partial^*E\cap\{z\}\) has \(\Hn\)-measure zero.  Since
\(|X_\eps|\le1\), dominated convergence applies there as well.  Letting
\(\eps\downarrow0\) in \eqref{eq:cutoff-GG} proves
\eqref{eq:radial-GG}.

Finally,
\[
   \frac12|\nu_E-e_z|^2
   =\frac12\bigl(|\nu_E|^2+|e_z|^2-2\nu_E\cdot e_z\bigr)
   =1-\nu_E\cdot e_z
\]
for \(\Hn\)-almost every \(x\in\partial^*E\).  Integration and
\eqref{eq:radial-GG} give \eqref{eq:osc-potential}.
\end{proof}

\begin{remark}[Unbounded sets]\label{rem:unbounded}
The identity extends to finite-measure finite-perimeter sets by introducing an
outer cut-off and passing to infinity.  The applications in this manuscript
either concern bounded sets or use the FJ theorem as quoted, so the bounded
version above is sufficient for every calculation made here.
\end{remark}

\subsection{Scaling and the large-deficit branch}

\begin{lemma}[Scaling]\label{lem:scaling}
Let \(\lambda>0\), \(a\in\R^n\), and \(F=a+\lambda E\).  If
\(|E|=|B_\rho|\), then \(|F|=|B_{\lambda\rho}|\) and
\[
   \Per(F)-\Per(B_{\lambda\rho})
      =\lambda^{n-1}\bigl(\Per(E)-\Per(B_\rho)\bigr),
\tag{2.5}\label{eq:scale-D}
\]
\[
   \beta(F)^2=\lambda^{n-1}\beta(E)^2.
\tag{2.6}\label{eq:scale-beta}
\]
\end{lemma}

\begin{proof}
Perimeter scales as \(\lambda^{n-1}\).  Also the map
\(x=a+\lambda y\) sends \(\partial^*E\) to \(\partial^*F\), preserves the
unit normal, and sends \(e_z(y)\) to \(e_{a+\lambda z}(x)\).  Therefore
\[
   \osc_{a+\lambda z}(F)=\lambda^{n-1}\osc_z(E).
\]
Taking the infimum over \(z\) proves \eqref{eq:scale-beta}.
\end{proof}

\begin{lemma}[Only small deficits require a subtle proof]\label{lem:large}
If \(|E|=|B_\rho|\) and
\(\Def(E)\ge\eta\Per(B_\rho)\) for some \(\eta>0\), then
\[
   \beta(E)^2\le 2(1+\eta^{-1})\Def(E).
\tag{2.7}\label{eq:large-branch}
\]
\end{lemma}

\begin{proof}
For any \(z\), \(|\nu_E-e_z|^2\le4\), and thus
\[
   \beta(E)^2\le\osc_z(E)\le2\Per(E).
\]
Since \(\Per(E)=\Per(B_\rho)+\Def(E)\) and
\(\Per(B_\rho)\le\eta^{-1}\Def(E)\), we obtain
\[
   \beta(E)^2\le2\Per(E)
   \le2(1+\eta^{-1})\Def(E).
\]
\end{proof}

Lemma~\ref{lem:large} explains the architecture of every effective proof:
one only has to convert a sufficiently small perimeter deficit into a
geometric regime in which the normal oscillation can be estimated explicitly.

\section{Attempt I: explicit nearly-spherical sets}\label{sec:graphs}

This attempt is complete.  We assume from the start that the set is a small
radial graph over a ball and prove a same-centre strong estimate by direct
calculation.  The balance assumptions remove the constant and translation
modes.  Their presence is not an inconvenience of the proof: a translation
of a ball has zero deficit and cannot be controlled as a graph about an
incorrect centre.

\subsection{Radial graphs, volume, and barycentre}

Let \(S=S^{n-1}\), with surface measure \(d\sigma\).  Recall
\[
   |S|=n\omega_n.
\tag{3.1}\label{eq:sphere-measure}
\]
Let \(u\in W^{1,\infty}(S)\) satisfy \(1+u>0\), and define
\[
   E_u:=\{r\theta:\theta\in S,\ 0\le r<1+u(\theta)\}.
\tag{3.2}\label{eq:radial-set}
\]
Throughout this section the norm is understood in the concrete form
\[
   \norm{u}_{W^{1,\infty}(S)}
   :=\max\{\norm{u}_{L^\infty(S)},
           \norm{\nabla_\tau u}_{L^\infty(S)}\}.
\]
Thus \(\norm{u}_{W^{1,\infty}}\le\delta\) separately implies
\(|u|\le\delta\) and \(|\nabla_\tau u|\le\delta\) almost everywhere.
For brevity write
\[
   a(\theta):=1+u(\theta),\qquad
   s(\theta):=|\nabla_\tau u(\theta)|,
\]
\[
   A:=\int_Ss^2\,d\sigma,\qquad
   B:=\int_Su^2\,d\sigma.
\tag{3.3}\label{eq:AB}
\]

\begin{lemma}[Volume and barycentre formulae]\label{lem:volume-bary}
For \(E_u\) defined in \eqref{eq:radial-set},
\[
   |E_u|=\frac1n\int_Sa^n\,d\sigma,
\tag{3.4}\label{eq:volume-graph}
\]
and
\[
   \int_{E_u}x\,dx=\frac1{n+1}\int_S\theta a^{n+1}\,d\sigma.
\tag{3.5}\label{eq:bary-graph}
\]
In particular, \(E_u\) has the volume and barycentre of \(B_1\) precisely when
\[
   \int_S(a^n-1)\,d\sigma=0,
   \qquad
   \int_S\theta a^{n+1}\,d\sigma=0.
\tag{3.6}\label{eq:balance}
\]
\end{lemma}

\begin{proof}
Use polar coordinates \(x=r\theta\), \(dx=r^{n-1}\,dr\,d\sigma(\theta)\):
\[
\begin{aligned}
   |E_u|
   &=\int_S\int_0^{a(\theta)}r^{n-1}\,dr\,d\sigma(\theta)  \\
   &=\frac1n\int_Sa(\theta)^n\,d\sigma(\theta).
\end{aligned}
\]
Similarly,
\[
\begin{aligned}
   \int_{E_u}x\,dx
   &=\int_S\int_0^{a(\theta)}r\theta\,r^{n-1}\,dr\,d\sigma(\theta) \\
   &=\frac1{n+1}\int_S\theta a(\theta)^{n+1}\,d\sigma(\theta).
\end{aligned}
\]
For \(u=0\), the two right sides equal \(|B_1|\) and \(0\), respectively,
because \(\int_S\theta\,d\sigma=0\).  This proves \eqref{eq:balance}.
\end{proof}

\subsection{Perimeter and fixed-centre normal oscillation}

\begin{lemma}[Geometry of the graph]\label{lem:graph-geometry}
For \(\Hn\)-almost every \(\theta\in S\), the point
\(x=a(\theta)\theta\) of \(\partial E_u\) has outer normal
\[
   \nu_{E_u}(a\theta)
   =\frac{a\theta-\nabla_\tau u}{\sqrt{a^2+s^2}},
\tag{3.7}\label{eq:graph-normal}
\]
and surface Jacobian
\[
   J_u(\theta)=a^{n-2}\sqrt{a^2+s^2}.
\tag{3.8}\label{eq:graph-jac}
\]
Consequently,
\[
   \Per(E_u)=\int_Sa^{n-2}\sqrt{a^2+s^2}\,d\sigma,
\tag{3.9}\label{eq:graph-perimeter}
\]
and the normal oscillation about the graph centre \(0\) is
\[
\begin{aligned}
   \osc_0(E_u)
   &=\int_Sa^{n-2}\bigl(\sqrt{a^2+s^2}-a\bigr)\,d\sigma \\
   &=\int_S
       \frac{a^{n-3}s^2}{\sqrt{1+s^2/a^2}+1}\,d\sigma.
\end{aligned}
\tag{3.10}\label{eq:graph-osc}
\]
\end{lemma}

\begin{proof}
Parametrise the boundary by
\[
   \Phi:S\to\R^n,\qquad \Phi(\theta)=a(\theta)\theta.
\]
If \(\tau\in T_\theta S\), then
\[
   D_\tau\Phi=a\tau+(D_\tau u)\theta.
\]
The vector \(a\theta-\nabla_\tau u\) is orthogonal to every such tangent
vector, since
\[
   (a\theta-\nabla_\tau u)\cdot
   \bigl(a\tau+(D_\tau u)\theta\bigr)
   =aD_\tau u-aD_\tau u=0.
\]
It points outward and has length \(\sqrt{a^2+s^2}\), proving
\eqref{eq:graph-normal}.

In an orthonormal tangent basis \(\tau_1,\ldots,\tau_{n-1}\), the Gram matrix
of \(D\Phi\) has entries
\[
   D_{\tau_i}\Phi\cdot D_{\tau_j}\Phi
   =a^2\delta_{ij}+D_{\tau_i}uD_{\tau_j}u.
\]
The determinant identity
\(\det(a^2I+v\otimes v)=a^{2(n-2)}(a^2+|v|^2)\) gives
\eqref{eq:graph-jac} and hence \eqref{eq:graph-perimeter}.

At \(x=a\theta\), the vector \(e_0(x)\) equals \(\theta\).  By
\eqref{eq:graph-normal},
\[
   1-\nu_{E_u}\cdot\theta
   =1-\frac{a}{\sqrt{a^2+s^2}}.
\]
Multiplying by \(J_u=a^{n-2}\sqrt{a^2+s^2}\), and recalling that
\(\osc_0=\int_{\partial E_u}(1-\nu\cdot e_0)\,d\Hn\), yields
\[
   \osc_0(E_u)
   =\int_Sa^{n-2}\bigl(\sqrt{a^2+s^2}-a\bigr)\,d\sigma.
\]
Finally,
\[
   \sqrt{a^2+s^2}-a
   =a\left(\sqrt{1+\frac{s^2}{a^2}}-1\right)
   =\frac{s^2/a}{\sqrt{1+s^2/a^2}+1},
\]
which proves the second line of \eqref{eq:graph-osc}.
\end{proof}

\subsection{Explicit constants}

For \(0<\delta<1\) and any real exponent \(p\), set
\[
   m_p(\delta):=\min_{t\in[1-\delta,1+\delta]}t^p,\qquad
   M_p(\delta):=\max_{t\in[1-\delta,1+\delta]}t^p.
\tag{3.11}\label{eq:mM}
\]
Define
\[
   L_n(\delta):=
   \frac{m_{n-3}(\delta)}
        {1+\sqrt{1+\bigl(\delta/(1-\delta)\bigr)^2}},
   \qquad
   U_n(\delta):=\frac12M_{n-3}(\delta),
\tag{3.12}\label{eq:LU}
\]
\[
   R_n(\delta):=
   \frac{n-1}{2}M_{n-3}(\delta)\bigl(1+(n-1)\delta\bigr),
\tag{3.13}\label{eq:R}
\]
\[
   C_0(\delta):=\frac{n-1}{2}M_{n-2}(\delta),\qquad
   C_1(\delta):=\frac n2M_{n-1}(\delta),
\tag{3.14}\label{eq:C01}
\]
\[
   \eta_n(\delta):=
   \delta^2\bigl(C_0(\delta)^2+nC_1(\delta)^2\bigr),
\tag{3.15}\label{eq:eta}
\]
and
\[
   d_n(\delta):=
   L_n(\delta)-\frac{R_n(\delta)}
   {2n\bigl(1-\eta_n(\delta)\bigr)}.
\tag{3.16}\label{eq:dn}
\]

\begin{theorem}[Explicit balanced nearly-spherical strong form]
\label{thm:graph-main}
Let \(n\ge2\).  Let \(E_u\) be defined by \eqref{eq:radial-set}, assume
\[
   |E_u|=|B_1|,\qquad \int_{E_u}x\,dx=0,
\tag{3.17}\label{eq:graph-balanced}
\]
and suppose
\[
   \norm{u}_{W^{1,\infty}(S)}\le\delta<1,\qquad
   \eta_n(\delta)<1,\qquad d_n(\delta)>0.
\tag{3.18}\label{eq:graph-hyp}
\]
Then
\[
   \beta(E_u)^2
   \le \frac{U_n(\delta)}{d_n(\delta)}
      \bigl(\Per(E_u)-\Per(B_1)\bigr).
\tag{3.19}\label{eq:graph-beta}
\]
More strongly, the same centre \(0\) satisfies
\[
   \int_{\partial^*E_u}
       \left|\nu_{E_u}-\frac{x}{|x|}\right|^2\,d\Hn
   \le
   \frac{2U_n(\delta)}{d_n(\delta)}
      \bigl(\Per(E_u)-\Per(B_1)\bigr),
\tag{3.20}\label{eq:graph-normal-bound}
\]
and
\[
   |E_u\mathbin{\triangle}B_1|^2
   \le
   \frac{M_{n-1}(\delta)^2\omega_n}
        {2(1-\eta_n(\delta))d_n(\delta)}
      \bigl(\Per(E_u)-\Per(B_1)\bigr).
\tag{3.21}\label{eq:graph-volume-bound}
\]
\end{theorem}

\begin{proof}
We break the proof into five explicit steps.

\smallskip\noindent\emph{Step 1: an upper bound for the fixed-centre
oscillation.}
By \eqref{eq:graph-osc} and the elementary inequality
\[
   \sqrt{1+t}+1\ge2\qquad(t\ge0),
\]
we have
\[
\begin{aligned}
   \osc_0(E_u)
   &=\int_S
       \frac{a^{n-3}s^2}{\sqrt{1+s^2/a^2}+1}\,d\sigma \\
   &\le \frac12M_{n-3}(\delta)\int_Ss^2\,d\sigma
    =U_n(\delta)A.
\end{aligned}
\tag{3.22}\label{eq:O-upper}
\]
Since \(\beta(E_u)^2=\inf_z\osc_z(E_u)\le\osc_0(E_u)\), it will suffice to
bound \(A\) by the perimeter deficit.

\smallskip\noindent\emph{Step 2: separate the perimeter deficit.}
Using \eqref{eq:graph-perimeter},
\[
\begin{aligned}
   \Per(E_u)-\Per(B_1)
   &=\int_S\bigl(a^{n-1}-1\bigr)\,d\sigma \\
   &\quad+
     \int_Sa^{n-2}\bigl(\sqrt{a^2+s^2}-a\bigr)\,d\sigma \\
   &=:I+G.
\end{aligned}
\tag{3.23}\label{eq:def-split}
\]
Since \(|u|\le\delta\) and \(s\le\delta\),
\[
   1-\delta\le a\le1+\delta,\qquad
   0\le\frac{s}{a}\le\frac{\delta}{1-\delta}.
\]
The exact identity
\[
   \sqrt{1+t}-1=\frac{t}{\sqrt{1+t}+1}
\]
therefore gives
\[
\begin{aligned}
   G
   &=\int_S
       \frac{a^{n-3}s^2}{\sqrt{1+s^2/a^2}+1}\,d\sigma \\
   &\ge
      \frac{m_{n-3}(\delta)}
      {1+\sqrt{1+\bigl(\delta/(1-\delta)\bigr)^2}}
      \int_Ss^2\,d\sigma \\
   &=L_n(\delta)A.
\end{aligned}
\tag{3.24}\label{eq:G-lower}
\]

\smallskip\noindent\emph{Step 3: estimate the part \(I\) by \(B\).}
The volume constraint in \eqref{eq:balance} yields
\[
   \int_S(a^n-1)\,d\sigma=0.
\]
Thus
\[
   I=\int_S f(u)\,d\sigma,
\qquad
   f(q):=(1+q)^{n-1}-1-\frac{n-1}{n}\bigl((1+q)^n-1\bigr).
\tag{3.25}\label{eq:def-f}
\]
Direct differentiation gives
\[
   f(0)=0,\qquad f'(0)=0,
\]
and
\[
\begin{aligned}
   f''(q)
   &=(n-1)(n-2)(1+q)^{n-3}
      -(n-1)^2(1+q)^{n-2} \\
   &=-(n-1)(1+q)^{n-3}\bigl(1+(n-1)q\bigr).
\end{aligned}
\tag{3.26}\label{eq:f-second}
\]
Taylor's formula with integral remainder reads
\[
   f(q)=q^2\int_0^1(1-t)f''(tq)\,dt.
\]
If \(|q|\le\delta\), then
\[
   |f''(tq)|
   \le(n-1)M_{n-3}(\delta)\bigl(1+(n-1)\delta\bigr).
\]
Since \(\int_0^1(1-t)\,dt=1/2\), we get
\[
   f(q)\ge-R_n(\delta)q^2.
\]
Integrating this inequality with \(q=u(\theta)\) proves
\[
   I\ge-R_n(\delta)B.
\tag{3.27}\label{eq:I-lower}
\]

\smallskip\noindent\emph{Step 4: remove the zero and first spherical
harmonic modes.}
The first balance identity in \eqref{eq:balance} can be written as
\[
   0=n\int_Su\,d\sigma+
     \int_S\bigl(a^n-1-nu\bigr)\,d\sigma.
\]
For \(|q|\le\delta\), Taylor's formula gives
\[
   |(1+q)^n-1-nq|
   \le\frac{n(n-1)}2M_{n-2}(\delta)q^2.
\]
After division by \(n\), it follows that
\[
   \left|\int_Su\,d\sigma\right|
   \le C_0(\delta)B.
\tag{3.28}\label{eq:mean-small}
\]
Likewise the barycentre identity gives
\[
   0=(n+1)\int_S\theta u\,d\sigma+
      \int_S\theta\bigl(a^{n+1}-1-(n+1)u\bigr)\,d\sigma.
\]
The estimate
\[
   |(1+q)^{n+1}-1-(n+1)q|
   \le\frac{n(n+1)}2M_{n-1}(\delta)q^2
\]
and \(|\theta|=1\) imply
\[
   \left|\int_S\theta u\,d\sigma\right|
   \le C_1(\delta)B.
\tag{3.29}\label{eq:first-small}
\]

Let \(u_0\) be the orthogonal projection of \(u\) onto the constants and let
\(u_1\) be its projection onto the first spherical harmonics.  Since
\(|S|=n\omega_n\),
\[
   \norm{u_0}_{L^2(S)}^2
   =\frac1{n\omega_n}\left(\int_Su\,d\sigma\right)^2.
\tag{3.30}\label{eq:proj0}
\]
The functions \(\theta_i/\sqrt{\omega_n}\), \(i=1,\ldots,n\), form an
orthonormal basis of the first harmonic space, because
\[
   \int_S\theta_i\theta_j\,d\sigma=\omega_n\delta_{ij}.
\]
Hence
\[
   \norm{u_1}_{L^2(S)}^2
   =\frac1{\omega_n}\left|\int_S\theta u\,d\sigma\right|^2.
\tag{3.31}\label{eq:proj1}
\]
From \eqref{eq:mean-small}, \eqref{eq:first-small},
\eqref{eq:proj0}, and \eqref{eq:proj1},
\[
\begin{aligned}
   \norm{u_0}_{2}^2+\norm{u_1}_{2}^2
   &\le
   \left(\frac{C_0(\delta)^2}{n\omega_n}
       +\frac{C_1(\delta)^2}{\omega_n}\right)B^2.
\end{aligned}
\tag{3.32}\label{eq:lowmodes1}
\]
Since \(|u|\le\delta\), we have
\[
   B=\int_Su^2\,d\sigma\le\delta^2|S|=\delta^2n\omega_n.
\]
Substitution in \eqref{eq:lowmodes1} gives
\[
   \norm{u_0}_{2}^2+\norm{u_1}_{2}^2
   \le\eta_n(\delta)B.
\tag{3.33}\label{eq:lowmodes2}
\]

The eigenvalues of \(-\Delta_S\) on spherical harmonics of degree \(k\) are
\[
   \lambda_k=k(k+n-2).
\]
Thus every harmonic of degree \(k\ge2\) has eigenvalue at least
\(\lambda_2=2n\).  The gradient energy has no contribution from the constant
mode, has a nonnegative contribution from the degree-one mode, and therefore
\[
\begin{aligned}
   A=\int_S|\nabla_\tau u|^2\,d\sigma
   &\ge2n\bigl(B-\norm{u_0}_2^2-\norm{u_1}_2^2\bigr) \\
   &\ge2n(1-\eta_n(\delta))B.
\end{aligned}
\tag{3.34}\label{eq:Poincare-balanced}
\]
Equivalently,
\[
   B\le\frac{A}{2n(1-\eta_n(\delta))}.
\tag{3.35}\label{eq:B-by-A}
\]

\smallskip\noindent\emph{Step 5: conclude.}
Putting \eqref{eq:def-split}, \eqref{eq:G-lower},
\eqref{eq:I-lower}, and \eqref{eq:B-by-A} together gives
\[
\begin{aligned}
   \Per(E_u)-\Per(B_1)
   &\ge L_n(\delta)A-R_n(\delta)B \\
   &\ge
   \left[
     L_n(\delta)-
     \frac{R_n(\delta)}{2n(1-\eta_n(\delta))}
   \right]A \\
   &=d_n(\delta)A.
\end{aligned}
\tag{3.36}\label{eq:D-by-A}
\]
Therefore \eqref{eq:O-upper} implies both
\[
   \beta(E_u)^2\le \osc_0(E_u)
   \le\frac{U_n(\delta)}{d_n(\delta)}
   \bigl(\Per(E_u)-\Per(B_1)\bigr)
\]
and, since the normal integral is \(2\osc_0(E_u)\),
\eqref{eq:graph-normal-bound}.

It remains only to prove \eqref{eq:graph-volume-bound}.  Along each ray from
the origin, \(E_u\mathbin{\triangle}B_1\) has one-dimensional radial measure
corresponding to the difference between \(a^n\) and \(1\).  Consequently,
\[
   |E_u\mathbin{\triangle}B_1|
   =\frac1n\int_S|a^n-1|\,d\sigma.
\tag{3.37}\label{eq:symdiff-formula}
\]
By the mean-value theorem,
\[
   |a^n-1|
   =|(1+u)^n-1|
   \le nM_{n-1}(\delta)|u|.
\]
Using Cauchy--Schwarz and \(|S|=n\omega_n\), we obtain
\[
\begin{aligned}
   |E_u\mathbin{\triangle}B_1|^2
   &\le M_{n-1}(\delta)^2
       \left(\int_S|u|\,d\sigma\right)^2 \\
   &\le M_{n-1}(\delta)^2(n\omega_n)B.
\end{aligned}
\tag{3.38}\label{eq:symdiff-B}
\]
Applying \eqref{eq:B-by-A} and \eqref{eq:D-by-A} to the final factor gives
\[
   |E_u\mathbin{\triangle}B_1|^2
   \le
   \frac{M_{n-1}(\delta)^2\omega_n}
        {2(1-\eta_n(\delta))d_n(\delta)}
       \bigl(\Per(E_u)-\Per(B_1)\bigr).
\]
This is \eqref{eq:graph-volume-bound}.
\end{proof}

\begin{corollary}[A simple displayed constant]\label{cor:graph-simple}
For every \(n\ge2\), set
\[
   \delta_*(n):=\frac1{100n^3}.
\tag{3.39}\label{eq:delta-star}
\]
If the hypotheses of Theorem~\ref{thm:graph-main} hold and
\(\norm{u}_{W^{1,\infty}(S)}\le\delta_*(n)\), then
\[
   \beta(E_u)^2
   \le3\bigl(\Per(E_u)-\Per(B_1)\bigr),
\tag{3.40}\label{eq:simple-beta}
\]
\[
   \int_{\partial^*E_u}
       \left|\nu_{E_u}-\frac{x}{|x|}\right|^2\,d\Hn
   \le6\bigl(\Per(E_u)-\Per(B_1)\bigr),
\tag{3.41}\label{eq:simple-normal}
\]
and
\[
   |E_u\mathbin{\triangle}B_1|^2
   \le3\omega_n\bigl(\Per(E_u)-\Per(B_1)\bigr).
\tag{3.42}\label{eq:simple-volume}
\]
\end{corollary}

\begin{proof}
We verify the numerical estimates without suppressing the elementary
calculation.  Let \(\delta=\delta_*(n)\).  Then
\[
   0<\delta\le\frac1{800}.
\tag{3.43}\label{eq:delta-bound}
\]
For the exponents \(p=n-3,n-2,n-1\), if \(p\ge0\), then
\[
   (1+\delta)^p\le e^{p\delta}
   \le e^{1/(100n^2)}
   \le e^{1/400}<1.004.
\]
The only negative exponent occurring is \(p=n-3=-1\) when \(n=2\); in that
case
\[
   M_{-1}(\delta)=\frac1{1-\delta}
   \le\frac{800}{799}<1.004.
\]
Thus
\[
   M_{n-3}(\delta),\ M_{n-2}(\delta),\ M_{n-1}(\delta)<1.004.
\tag{3.44}\label{eq:Mnumeric}
\]
Similarly, if \(n\ge4\), Bernoulli's inequality gives
\[
   (1-\delta)^{n-3}\ge1-(n-3)\delta
   \ge1-\frac1{100n^2}>0.997.
\]
For \(n=3\), \(m_{n-3}=m_0=1\); for \(n=2\),
\[
   m_{-1}(\delta)=\frac1{1+\delta}>\frac{800}{801}>0.997.
\]
Hence
\[
   m_{n-3}(\delta)>0.997.
\tag{3.45}\label{eq:mnumeric}
\]
Furthermore,
\[
   \frac{\delta}{1-\delta}\le\frac1{799},
\]
and so
\[
   1+\sqrt{1+\left(\frac{\delta}{1-\delta}\right)^2}
   <2.000002.
\]
Using \eqref{eq:mnumeric}, we conclude that
\[
   L_n(\delta)>\frac{0.997}{2.000002}>0.498,
   \qquad
   U_n(\delta)<\frac{1.004}{2}=0.502.
\tag{3.46}\label{eq:LUnumeric}
\]

From \eqref{eq:C01} and \eqref{eq:Mnumeric},
\[
   C_0(\delta)<0.502n,\qquad C_1(\delta)<0.502n.
\]
Therefore
\[
\begin{aligned}
   \eta_n(\delta)
   &<\delta^2(0.502^2n^2+0.502^2n^3) \\
   &\le0.504008\,\delta^2n^3
    =\frac{0.504008}{10000n^3}
    <10^{-4}.
\end{aligned}
\tag{3.47}\label{eq:etanumeric}
\]
Also
\[
   1+(n-1)\delta
   \le1+\frac1{100n^2}
   \le1.0025,
\]
and hence, using \eqref{eq:Mnumeric} and \eqref{eq:etanumeric},
\[
\begin{aligned}
   \frac{R_n(\delta)}{2n(1-\eta_n(\delta))}
   &<
   \frac{n-1}{4n}
   \frac{1.004\cdot1.0025}{1-10^{-4}} \\
   &<0.252.
\end{aligned}
\tag{3.48}\label{eq:Rnumeric}
\]
Equations \eqref{eq:LUnumeric} and \eqref{eq:Rnumeric} yield
\[
   d_n(\delta)>0.498-0.252=0.246>0.2.
\tag{3.49}\label{eq:dnumeric}
\]
It follows that
\[
   \frac{U_n(\delta)}{d_n(\delta)}
   <\frac{0.502}{0.246}<3,
\]
proving \eqref{eq:simple-beta} and \eqref{eq:simple-normal}.
Finally,
\[
   \frac{M_{n-1}(\delta)^2}
        {2(1-\eta_n(\delta))d_n(\delta)}
   <
   \frac{1.004^2}{2(1-10^{-4})0.246}<3,
\]
which proves \eqref{eq:simple-volume}.
\end{proof}

\begin{corollary}[Translation and scaling]\label{cor:scaled-graph}
Let \(\rho>0\), \(z\in\R^n\), and suppose
\[
   E=z+\rho E_u,
\]
where \(E_u\) satisfies the hypotheses of Corollary~\ref{cor:graph-simple}.
Then the same centre \(z\) satisfies
\[
   \rho^{-n-1}|E\mathbin{\triangle}B_\rho(z)|^2
   +
   \int_{\partial^*E}
      \left|\nu_E-\frac{x-z}{|x-z|}\right|^2\,d\Hn
   \le
   (3\omega_n+6)\bigl(\Per(E)-\Per(B_\rho)\bigr).
\tag{3.50}\label{eq:scaled-same-centre}
\]
\end{corollary}

\begin{proof}
Under \(x=z+\rho y\), symmetric-difference measure scales by \(\rho^n\),
while perimeter and the boundary normal integral scale by \(\rho^{n-1}\).
Therefore \eqref{eq:simple-volume} becomes
\[
   |E\mathbin{\triangle}B_\rho(z)|^2
   \le3\omega_n\rho^{n+1}
        \bigl(\Per(E)-\Per(B_\rho)\bigr),
\]
and \eqref{eq:simple-normal} becomes
\[
   \int_{\partial^*E}
      \left|\nu_E-\frac{x-z}{|x-z|}\right|^2\,d\Hn
   \le6\bigl(\Per(E)-\Per(B_\rho)\bigr).
\]
Adding the two inequalities proves \eqref{eq:scaled-same-centre}.
\end{proof}

\subsection{The quadratic constant}

\begin{proposition}[Degree-two test]\label{prop:degree-two}
Let \(Y\not\equiv0\) be an even spherical harmonic of degree two on \(S\).
There exist constants \(c_\eps=O(\eps^2)\) such that
\[
   u_\eps=\eps Y+c_\eps,\qquad
   |E_{u_\eps}|=|B_1|,\qquad
   \int_{E_{u_\eps}}x\,dx=0.
\]
For this family,
\[
   \frac{\beta(E_{u_\eps})^2}
   {\Per(E_{u_\eps})-\Per(B_1)}
   \longrightarrow \frac{2n}{n+1}
   \qquad\text{as }\eps\to0.
\tag{3.51}\label{eq:sharp-constant}
\]
\end{proposition}

\begin{proof}
Because \(Y\) is a degree-two spherical harmonic,
\[
   \int_SY\,d\sigma=0,\qquad
   -\Delta_SY=2nY.
\tag{3.52}\label{eq:Yfacts}
\]
Since \(Y(-\theta)=Y(\theta)\), every \(u_\eps=\eps Y+c_\eps\) is even.
Therefore \(E_{u_\eps}\) is centrally symmetric, and its barycentre is zero.

Define
\[
   F(c,\eps):=\int_S(1+c+\eps Y)^n\,d\sigma-|S|.
\]
We have \(F(0,0)=0\) and
\[
   \partial_cF(0,0)=n|S|>0.
\]
The implicit function theorem gives a smooth function \(c_\eps\), for small
\(\eps\), such that \(F(c_\eps,\eps)=0\).  Expanding to second order,
\[
\begin{aligned}
   0
   &=n|S|c_\eps+
     \frac{n(n-1)}2\eps^2\int_SY^2\,d\sigma
     +O(\eps^3),
\end{aligned}
\]
so
\[
   c_\eps
   =-\frac{n-1}{2|S|}\eps^2\int_SY^2\,d\sigma+O(\eps^3).
\tag{3.53}\label{eq:ceps}
\]

The perimeter expansion proved directly from
\eqref{eq:graph-perimeter} and the volume constraint is
\[
   \Per(E_u)-\Per(B_1)
   =
   \frac12\int_S|\nabla_\tau u|^2\,d\sigma
   -\frac{n-1}{2}\int_Su^2\,d\sigma
   +O(\norm{u}_{W^{1,\infty}})\!
      \int_S(u^2+|\nabla_\tau u|^2)\,d\sigma.
\tag{3.54}\label{eq:D-secondvariation}
\]
For completeness, this formula follows by expanding
\[
   a^{n-1}=1+(n-1)u+\frac{(n-1)(n-2)}2u^2+O(|u|^3)
\]
and
\[
   \sqrt{1+\frac{s^2}{a^2}}
   =1+\frac12s^2+O\bigl(\norm{u}_{W^{1,\infty}}s^2\bigr),
\]
then using
\[
   \int_Su\,d\sigma
   =-\frac{n-1}{2}\int_Su^2\,d\sigma+O(\norm{u}_\infty)\int_Su^2\,d\sigma,
\]
which is the second-order form of the volume constraint.
Since \(c_\eps=O(\eps^2)\), \eqref{eq:Yfacts} and
\eqref{eq:D-secondvariation} give
\[
\begin{aligned}
   \Per(E_{u_\eps})-\Per(B_1)
   &=\frac{\eps^2}{2}
      \left(\int_S|\nabla_\tau Y|^2\,d\sigma
       -(n-1)\int_SY^2\,d\sigma\right)+o(\eps^2) \\
   &=\frac{n+1}{2}\eps^2\int_SY^2\,d\sigma+o(\eps^2).
\end{aligned}
\tag{3.55}\label{eq:D-Y}
\]

It remains to expand the infimum over centres in \(\beta\).
By \eqref{eq:graph-osc},
\[
   \osc_0(E_{u_\eps})
   =\frac{\eps^2}{2}\int_S|\nabla_\tau Y|^2\,d\sigma+o(\eps^2)
   =n\eps^2\int_SY^2\,d\sigma+o(\eps^2).
\tag{3.56}\label{eq:O0-Y}
\]
We now check that translating the comparison centre cannot decrease the
quadratic term.  If \(z=\eps b\) with \(b\in\R^n\) bounded, then, uniformly on
\(S\),
\[
   \nu_{E_{u_\eps}}=\theta-\eps\nabla_\tau Y+O(\eps^2),
\]
and
\[
   \frac{(1+u_\eps)\theta-\eps b}
        {|(1+u_\eps)\theta-\eps b|}
   =\theta-\eps\bigl(b-(b\cdot\theta)\theta\bigr)+O(\eps^2).
\]
The vector field
\[
   b-(b\cdot\theta)\theta=\nabla_\tau(b\cdot\theta)
\]
is the gradient of a degree-one spherical harmonic.  Orthogonality of
spherical harmonics gives
\[
   \int_S\nabla_\tau Y\cdot\nabla_\tau(b\cdot\theta)\,d\sigma
   =\int_S(-\Delta_SY)(b\cdot\theta)\,d\sigma=0.
\tag{3.57}\label{eq:orthog}
\]
Therefore
\[
\begin{aligned}
   \osc_{\eps b}(E_{u_\eps})
   &=
   \frac{\eps^2}{2}\int_S
   \left|\nabla_\tau Y-\nabla_\tau(b\cdot\theta)\right|^2\,d\sigma
   +o(\eps^2) \\
   &=
   \frac{\eps^2}{2}\int_S|\nabla_\tau Y|^2\,d\sigma
   +\frac{\eps^2}{2}\int_S|\nabla_\tau(b\cdot\theta)|^2\,d\sigma
   +o(\eps^2).
\end{aligned}
\tag{3.58}\label{eq:center-quadratic}
\]
The additional term is nonnegative and vanishes only for \(b=0\).
We now justify that it is enough to consider such bounded \(b\).  By
Proposition~\ref{prop:radial-GG},
\[
   \osc_z(E_{u_\eps})
   =\Per(E_{u_\eps})-(n-1)\int_{E_{u_\eps}}\frac{dx}{|x-z|}.
\]
For bounded sets the potential in this formula is continuous in \(z\) and
tends to \(0\) as \(|z|\to\infty\).  Hence the infimum defining
\(\beta(E_{u_\eps})\) is attained.  Let \(z_\eps\) be a minimising centre.
Since \(\osc_0(E_{u_\eps})=O(\eps^2)\), we have
\(\osc_{z_\eps}(E_{u_\eps})=O(\eps^2)\).  The convergence
\(E_{u_\eps}\to B_1\) in \(C^1\) then implies \(z_\eps\to0\): otherwise,
after passage to a subsequence, the limiting oscillation would be
\(\osc_z(B_1)>0\) for some \(z\ne0\), or \(\Per(B_1)>0\) if
\(|z_\eps|\to\infty\).

For \(|z|+|\eps|\) small, Taylor expansion of the two smooth boundary vector
fields used above is uniform in \(\theta\) and gives
\[
\begin{aligned}
   \osc_z(E_{u_\eps})
   &=
   \frac12\int_S
      \left|-\eps\nabla_\tau Y+
       \bigl(z-(z\cdot\theta)\theta\bigr)\right|^2\,d\sigma \\
   &\quad+O\bigl((|\eps|+|z|)^3\bigr).
\end{aligned}
\tag{3.59}\label{eq:center-general}
\]
The cross term is zero by \eqref{eq:orthog}, and
\[
\begin{aligned}
   \int_S\left|z-(z\cdot\theta)\theta\right|^2\,d\sigma
   &=|z|^2|S|-\int_S(z\cdot\theta)^2\,d\sigma \\
   &=(n-1)\omega_n|z|^2.
\end{aligned}
\tag{3.60}\label{eq:translation-coercive}
\]
If \(|z_\eps|/|\eps|\to\infty\), then
\eqref{eq:center-general}--\eqref{eq:translation-coercive} yield
\(\osc_{z_\eps}(E_{u_\eps})\gg\eps^2\), contradicting
\(\osc_{z_\eps}(E_{u_\eps})\le\osc_0(E_{u_\eps})=O(\eps^2)\).
Thus \(z_\eps=\eps b_\eps\) with \(b_\eps\) bounded.  Formula
\eqref{eq:center-quadratic} applies along subsequences and its nonnegative
translation term is minimised at \(b=0\).
Thus
\[
   \beta(E_{u_\eps})^2
   =n\eps^2\int_SY^2\,d\sigma+o(\eps^2).
\tag{3.61}\label{eq:beta-Y}
\]
Dividing \eqref{eq:beta-Y} by \eqref{eq:D-Y} proves
\eqref{eq:sharp-constant}.
\end{proof}

\begin{remark}[Meaning of the constant]\label{rem:constant}
Theorem~\ref{thm:graph-main} is intentionally conservative: the simple
constant \(3\) is easy to verify in every dimension.  Proposition
\ref{prop:degree-two} identifies the correct leading-order obstruction for
the oscillation term.  In particular, a claimed asymptotic constant for
\(\beta(E)^2/\Def(E)\) which is smaller than \(2n/(n+1)\) is incompatible with
degree-two radial perturbations.
\end{remark}

\section{Attempt II: quantitative selection and regularisation}
\label{sec:selection}

We begin with a decomposition-based selection attempt.  Its advantage is
that all quotient-preservation estimates can be obtained immediately.  Its
defect, made explicit in Remark~\ref{rem:M-sharp}, is that the selected set
does not immediately enter the additive perimeter-almost-minimiser class
needed for regularity.  The corrected regularising procedure appears in
Subsection~\ref{ssec:true-beta-selection}; it returns to the Fusco--Julin
choice of penalising \(\beta^2\) itself, but quantifies the intermediate
annular-separation step which makes that choice regularising.

For the first attempt, use the exact positive decomposition
\[
   \beta(E)^2=\Def(E)+\mathcal C(E),
\tag{4.1}\label{eq:beta-decomposition-selection}
\]
where \(\mathcal C\) is a potential deficit.  It is the analogue of the
smooth asymmetry preserved by the single-set selection map in Plan~1, but
only at the quotient-preservation level.

\subsection{The potential deficit}

For every measurable \(U\subset\R^n\) with finite measure, define
\[
   r_U:=\left(\frac{|U|}{\omega_n}\right)^{1/n},
   \qquad
   p(U):=\Per(B_{r_U})
          =n\omega_n^{1/n}|U|^{(n-1)/n},
\tag{4.2}\label{eq:pU}
\]
\[
   M(U):=\sup_{z\in\R^n}\int_U\frac{dx}{|x-z|},
\tag{4.3}\label{eq:MU}
\]
and
\[
   \Def(U):=\Per(U)-p(U),\qquad
   \mathcal C(U):=p(U)-(n-1)M(U).
\tag{4.4}\label{eq:DC-general}
\]
For the empty set, all four quantities are set equal to zero.

\begin{lemma}[Positivity and decomposition]\label{lem:C-decomposition}
For every finite-perimeter set \(U\) of finite measure,
\[
   \Def(U)\ge0,\qquad \mathcal C(U)\ge0,
\tag{4.5}\label{eq:DC-positive}
\]
and
\[
   \beta(U)^2=\Def(U)+\mathcal C(U),
\tag{4.6}\label{eq:beta-DC}
\]
where the radius in the definition of \(\beta(U)\) is \(r_U\).
Under a dilation \(U\mapsto a+\lambda U\),
\[
   \Def(a+\lambda U)=\lambda^{n-1}\Def(U),\qquad
   \mathcal C(a+\lambda U)=\lambda^{n-1}\mathcal C(U).
\tag{4.7}\label{eq:DC-scale}
\]
\end{lemma}

\begin{proof}
The isoperimetric inequality gives
\(\Per(U)\ge\Per(B_{r_U})=p(U)\), proving \(\Def(U)\ge0\).
For each \(z\), the radially decreasing rearrangement inequality applied to
\(x\mapsto |x-z|^{-1}\) gives
\[
   \int_U\frac{dx}{|x-z|}
   \le\int_{B_{r_U}(z)}\frac{dx}{|x-z|}
   =n\omega_n\int_0^{r_U}s^{n-2}\,ds
   =\frac{p(U)}{n-1}.
\]
Taking the supremum in \(z\) proves \(\mathcal C(U)\ge0\).

By Proposition~\ref{prop:radial-GG},
\[
\begin{aligned}
   \beta(U)^2
   &=\Per(U)-(n-1)M(U) \\
   &=\bigl(\Per(U)-p(U)\bigr)
     +\bigl(p(U)-(n-1)M(U)\bigr) \\
   &=\Def(U)+\mathcal C(U).
\end{aligned}
\]
Finally, perimeter and \(p\) scale by \(\lambda^{n-1}\); the change of
variables \(x=a+\lambda y\) in \eqref{eq:MU} gives
\(M(a+\lambda U)=\lambda^{n-1}M(U)\).  This proves
\eqref{eq:DC-scale}.
\end{proof}

\begin{remark}[The new badness ratio]\label{rem:QC}
If \(\mathcal C(E)=0\), then \eqref{eq:beta-DC} gives
\(\beta(E)^2=\Def(E)\), and there is nothing to prove.  Otherwise define
\[
   Q_{\mathcal C}(E):=\frac{\Def(E)}{\mathcal C(E)}.
\tag{4.8}\label{eq:QC}
\]
An estimate \(Q_{\mathcal C}(E)\ge c>0\) is equivalent to
\[
   \beta(E)^2=\Def(E)+\mathcal C(E)
   \le(1+c^{-1})\Def(E).
\]
Thus \(Q_{\mathcal C}\) is a legitimate selection quotient for the FJ
problem.  It is also unchanged by translation and dilation, by
\eqref{eq:DC-scale}.
\end{remark}

\subsection{A continuity estimate for the preserved denominator}

The following estimate is the point at which the potential deficit is better
suited to selection than the full quantity \(\beta^2\): it is continuous in
measure, with an explicit modulus.

\begin{lemma}[Continuity of the maximal radial potential]\label{lem:M-cont}
Let \(U,V\subset\R^n\) have finite measure.  Then
\[
   |M(U)-M(V)|
   \le K_n\,|U\mathbin{\triangle}V|^{(n-1)/n},
   \qquad
   K_n:=\frac{n}{n-1}\omega_n^{1/n}.
\tag{4.9}\label{eq:M-cont}
\]
\end{lemma}

\begin{proof}
Fix \(z\in\R^n\) and set \(A=U\mathbin{\triangle}V\).  Since the integrand
is nonnegative,
\[
   \left|\int_U\frac{dx}{|x-z|}
          -\int_V\frac{dx}{|x-z|}\right|
   \le\int_A\frac{dx}{|x-z|}.
\]
Among all sets of measure \(|A|\), the final integral is maximised by the
ball centred at \(z\).  If \(|A|=\omega_n r^n\), then
\[
\begin{aligned}
   \int_A\frac{dx}{|x-z|}
   &\le\int_{B_r(z)}\frac{dx}{|x-z|} \\
   &=n\omega_n\int_0^r s^{n-2}\,ds
     =\frac{n}{n-1}\omega_n^{1/n}|A|^{(n-1)/n}.
\end{aligned}
\]
This bound is independent of \(z\).  Taking the supremum over \(z\) once
with \(U\) and once with \(V\) proves \eqref{eq:M-cont}.
\end{proof}

\begin{corollary}[Continuity of \(\mathcal C\) on bounded-volume classes]
\label{cor:C-cont}
Let \(0<m_-\le |U|,|V|\le m_+<\infty\).  Then
\[
   |\mathcal C(U)-\mathcal C(V)|
   \le L_p(m_-,m_+)|U\mathbin{\triangle}V|
      +n\omega_n^{1/n}|U\mathbin{\triangle}V|^{(n-1)/n},
\tag{4.10}\label{eq:C-cont}
\]
where one may take
\[
   L_p(m_-,m_+)
   :=(n-1)\omega_n^{1/n}m_-^{-1/n}.
\tag{4.11}\label{eq:Lp}
\]
\end{corollary}

\begin{proof}
The derivative of
\(m\mapsto n\omega_n^{1/n}m^{(n-1)/n}\) is
\((n-1)\omega_n^{1/n}m^{-1/n}\), bounded by \(L_p\) on
\([m_-,m_+]\).  Combine this Lipschitz estimate for \(p\) with
Lemma~\ref{lem:M-cont} and definition \eqref{eq:DC-general}.
\end{proof}

\begin{remark}[The singular modulus is sharp]\label{rem:M-sharp}
The exponent in Lemma~\ref{lem:M-cont} cannot be improved to a Lipschitz
bound in the symmetric-difference measure.  Indeed, for
\(U=B_r(0)\) and \(V=\varnothing\),
\[
   M(U)-M(V)
   =\frac{n\omega_n}{n-1}r^{n-1}
   =\frac{n}{n-1}\omega_n^{1/n}
      |U\mathbin{\triangle}V|^{(n-1)/n}.
\]
Thus local replacements in a ball of radius \(r\) naturally change the soft
constraint by order \(r^{n-1}\), not order \(r^n\).  This is why the
Plan~1 selection architecture transfers, but its lower-order regularity
error does not transfer automatically.
\end{remark}

\subsection{The single-set selection map}

We first work in a container \(B_R\), with \(R>1\).  This is analogous to
the bounded core of Plan~1.  For the direct-\(\beta\) regularisation used
below, the computable confinement step needed to pass to arbitrary sets is
proved in Lemma~\ref{lem:computable-confinement}.

Fix a volume-penalty coefficient \(\Lambda>0\).  Given an input set
\[
   E_0\subset B_R,\qquad |E_0|=|B_1|=\omega_n,\qquad
   \epsilon:=\mathcal C(E_0)>0,
\]
write
\[
   \delta:=\Def(E_0),\qquad q:=\frac{\delta}{\epsilon}.
\tag{4.12}\label{eq:input-q}
\]
For \(0<\tau<1\), define
\[
   h_{\epsilon,\tau}(s)
   :=\sqrt{\epsilon^2+\tau^2(s-\epsilon)^2}
   \qquad (s\ge0),
\tag{4.13}\label{eq:h}
\]
and minimise
\[
   \mathcal J_{E_0,\tau,\Lambda}(U)
   :=
   \Def(U)+\Lambda\bigl||U|-\omega_n\bigr|
   +h_{\epsilon,\tau}\bigl(\mathcal C(U)\bigr)
\tag{4.14}\label{eq:J-FJ}
\]
among finite-perimeter sets \(U\subset B_R\).  Notice that
\[
   \mathcal J_{E_0,\tau,\Lambda}(E_0)=\delta+\epsilon.
\tag{4.15}\label{eq:J-input}
\]

\begin{lemma}[Elementary facts about the soft constraint]\label{lem:h}
For all \(s,t\ge0\),
\[
   h_{\epsilon,\tau}(s)\ge\epsilon,
\qquad
   |h_{\epsilon,\tau}(s)-h_{\epsilon,\tau}(t)|
   \le\tau|s-t|.
\tag{4.16}\label{eq:h-properties}
\]
\end{lemma}

\begin{proof}
The first statement follows immediately from \eqref{eq:h}.  Differentiating
for \(s>0\),
\[
   h_{\epsilon,\tau}'(s)
   =
   \frac{\tau^2(s-\epsilon)}
        {\sqrt{\epsilon^2+\tau^2(s-\epsilon)^2}},
\]
and hence \(|h_{\epsilon,\tau}'(s)|\le\tau\).  The Lipschitz estimate
follows from the mean-value theorem.
\end{proof}

\begin{theorem}[Single-set FJ selection]\label{thm:FJ-selection}
Let \(E_0,\epsilon,\delta,q\) be as in \eqref{eq:input-q}, and assume
\[
   0<q\le\frac1{144},\qquad
   \delta<\frac{\Lambda\omega_n}{2}.
\tag{4.17}\label{eq:q-select-small}
\]
Choose
\[
   \tau:=q^{1/4}.
\tag{4.18}\label{eq:tau}
\]
Then \(\mathcal J_{E_0,\tau,\Lambda}\) has a minimiser \(F\subset B_R\).
Every minimiser satisfies
\[
   \Def(F)\le\delta,\qquad
   \bigl||F|-\omega_n\bigr|\le\frac{\delta}{\Lambda},
\tag{4.19}\label{eq:selected-def-volume}
\]
\[
   |\mathcal C(F)-\epsilon|
   \le\epsilon\sqrt3\,q^{1/4}
   \le\frac{\epsilon}{2},
\tag{4.20}\label{eq:selected-C}
\]
and hence
\[
   Q_{\mathcal C}(F)
   =\frac{\Def(F)}{\mathcal C(F)}
   \le2\,Q_{\mathcal C}(E_0)=2q.
\tag{4.21}\label{eq:selected-ratio}
\]
If
\[
   \lambda:=\left(\frac{\omega_n}{|F|}\right)^{1/n},
   \qquad \widetilde F:=\lambda F,
\tag{4.22}\label{eq:F-normalise}
\]
then \(|\widetilde F|=|B_1|\) and
\[
   Q_{\mathcal C}(\widetilde F)=Q_{\mathcal C}(F)\le2q.
\tag{4.23}\label{eq:selected-normal-ratio}
\]
\end{theorem}

\begin{proof}
\emph{Existence.}
Let \((U_j)\) be a minimising sequence.  By comparison with \(E_0\), we may
assume
\[
   \mathcal J_{E_0,\tau,\Lambda}(U_j)\le\delta+\epsilon+1.
\]
All three terms in \eqref{eq:J-FJ} are nonnegative by
Lemma~\ref{lem:C-decomposition} and Lemma~\ref{lem:h}.  Hence
\(\Def(U_j)\le\delta+1\).  Since \(U_j\subset B_R\),
\[
   \Per(U_j)=\Def(U_j)+p(U_j)
   \le \delta+1+\Per(B_R),
\]
so the perimeters are uniformly bounded.  Compactness in \(BV(B_R)\)
gives a subsequence, not relabelled, and a finite-perimeter set
\(F\subset B_R\) such that
\[
   {\bf1}_{U_j}\to{\bf1}_{F}\quad\text{in }L^1(B_R),
   \qquad
   \Per(F)\le\liminf_j\Per(U_j).
\]
The volume converges.  The function \(p(U)\) therefore converges, and
Lemma~\ref{lem:M-cont} gives \(M(U_j)\to M(F)\), hence
\(\mathcal C(U_j)\to\mathcal C(F)\).  It follows that
\(\mathcal J_{E_0,\tau,\Lambda}\) is lower semicontinuous along this
subsequence and \(F\) is a minimiser.

\emph{Quantitative preservation.}
Minimality against \(E_0\) and \eqref{eq:J-input} give
\[
   \Def(F)+\Lambda\bigl||F|-\omega_n\bigr|
   +h_{\epsilon,\tau}(\mathcal C(F))
   \le\delta+\epsilon.
\tag{4.24}\label{eq:minimal-compare}
\]
Since \(h_{\epsilon,\tau}\ge\epsilon\), dropping the two other nonnegative
terms in turn proves \eqref{eq:selected-def-volume}.  Dropping
\(\Def(F)\) and the volume term instead gives
\[
   h_{\epsilon,\tau}(\mathcal C(F))\le\epsilon+\delta.
\]
Squaring this inequality and using \(\delta=q\epsilon\), we obtain
\[
\begin{aligned}
   \tau^2(\mathcal C(F)-\epsilon)^2
   &\le2\epsilon\delta+\delta^2
     =\epsilon^2(2q+q^2), \\
   |\mathcal C(F)-\epsilon|
   &\le\epsilon\frac{\sqrt{2q+q^2}}{\tau}
     \le\epsilon\sqrt3\,q^{1/4},
\end{aligned}
\]
because \(q\le1\) and \(\tau=q^{1/4}\).  If \(q\le1/144\), then
\(\sqrt3q^{1/4}\le1/2\), proving \eqref{eq:selected-C}.
Consequently \(\mathcal C(F)\ge\epsilon/2\), and
\[
   \frac{\Def(F)}{\mathcal C(F)}
   \le\frac{\delta}{\epsilon/2}=2q.
\]

\emph{Normalisation.}
The bound on \(|F|\) in \eqref{eq:selected-def-volume}, together with
\eqref{eq:q-select-small}, implies \(|F|\ge\omega_n/2>0\); therefore
\eqref{eq:F-normalise} is well-defined.  Formula \eqref{eq:DC-scale}
shows that the two terms in \(Q_{\mathcal C}\) are multiplied by the same
factor \(\lambda^{n-1}\), proving \eqref{eq:selected-normal-ratio}.
\end{proof}

\begin{remark}[Difference from Plan~1]\label{rem:selection-difference}
Plan~1 preserves a smooth volume asymmetry while minimising torsional energy.
Here the selected energy is the perimeter deficit and the preserved
asymmetry is \(\mathcal C\), the non-perimeter part of the FJ oscillation.
This is not merely a change of notation.  Penalising \(\beta^2\) directly
would penalise a quantity containing \(\Per(U)\), obscuring the perimeter
coercivity needed for compactness and regularity.  Decomposition
\eqref{eq:beta-DC} removes that obstruction.
\end{remark}

\subsection{The selected set is an explicitly quantified near-minimiser}

The counterpart of the free-boundary quasi-minimality in Plan~1 is now a
perimeter near-minimality estimate.  It is an actual consequence of the
selection theorem, not an additional assumption.

\begin{proposition}[Local near-minimality]\label{prop:near-min}
Let \(F\) be selected by Theorem~\ref{thm:FJ-selection}.  Assume in addition
\(\delta\le\Lambda\omega_n/4\), so \(|F|\in[3\omega_n/4,5\omega_n/4]\).
Set
\[
   r_0:=\left(\frac{\omega_n}{4\omega_n}\right)^{1/n}=4^{-1/n},
   \qquad
   L_*:=(n-1)\omega_n^{1/n}\left(\frac{\omega_n}{2}\right)^{-1/n}.
\tag{4.25}\label{eq:r0-L}
\]
If \(0<r\le r_0\) and \(V\subset B_R\) is a finite-perimeter competitor
satisfying
\[
   F\mathbin{\triangle}V\Subset B_r(x)\Subset B_R,
\tag{4.26}\label{eq:local-competitor}
\]
then
\[
   \Per(F;B_r(x))
   \le\Per(V;B_r(x))
      +A_0\,r^n+A_1\,\tau r^{n-1},
\tag{4.27}\label{eq:near-min}
\]
where
\[
   A_0:=\omega_n\bigl(\Lambda+(1+\tau)L_*\bigr),
   \qquad
   A_1:=n\omega_n.
\tag{4.28}\label{eq:A01}
\]
\end{proposition}

\begin{proof}
Since \(F\mathbin{\triangle}V\subset B_r(x)\),
\[
   \bigl||V|-|F|\bigr|\le |F\mathbin{\triangle}V|
   \le\omega_n r^n\le\frac{\omega_n}{4}.
\]
Thus \(|V|\ge\omega_n/2\), and the derivative estimate in
\eqref{eq:Lp} gives
\[
   |p(F)-p(V)|\le L_*|F\mathbin{\triangle}V|.
\tag{4.29}\label{eq:p-local}
\]
Minimality of \(F\), cancellation of the perimeter outside \(B_r(x)\), and
Lemma~\ref{lem:h} imply
\[
\begin{aligned}
   \Per(F;B_r(x))
   &\le \Per(V;B_r(x))
      +|p(F)-p(V)|
      +\Lambda|F\mathbin{\triangle}V| \\
   &\quad+\tau|\mathcal C(F)-\mathcal C(V)|.
\end{aligned}
\tag{4.30}\label{eq:local-before-C}
\]
By Corollary~\ref{cor:C-cont} and \eqref{eq:p-local},
\[
   |\mathcal C(F)-\mathcal C(V)|
   \le L_*|F\mathbin{\triangle}V|
      +n\omega_n^{1/n}|F\mathbin{\triangle}V|^{(n-1)/n}.
\]
Insert this bound in \eqref{eq:local-before-C}, use
\(|F\mathbin{\triangle}V|\le\omega_n r^n\), and calculate
\[
   n\omega_n^{1/n}(\omega_n r^n)^{(n-1)/n}
   =n\omega_n r^{n-1}.
\]
The result is \eqref{eq:near-min}--\eqref{eq:A01}.
\end{proof}

\begin{remark}[The exact regularity burden]\label{rem:near-min-burden}
Estimate \eqref{eq:near-min} contains a small error of perimeter order,
\(\tau r^{n-1}\), in addition to a lower-order volume error.  A constructive
regularity argument must retain this small coefficient rather than hide it
inside an unnamed ``almost-minimiser'' theorem.  This is the FJ counterpart
of Plan~1's Alt--Caffarelli input.
\end{remark}

\begin{lemma}[Density and measure closeness give Hausdorff closeness]
\label{lem:FJ-density-Hausdorff}
Let \(G\subset\R^n\) be a bounded finite-perimeter set with
\(|G|=|B_1|\), represented so that its topological boundary agrees with
the support of its perimeter measure.  Assume there exist \(c_d>0\) and
\(0<r_d\le1/4\) such that
\[
   |G\cap B_r(x)|\ge c_d r^n,\qquad
   |B_r(x)\setminus G|\ge c_d r^n
\tag{4.31}\label{eq:density-input}
\]
for every \(x\in\partial G\) and every \(0<r\le r_d\).
Let \(z\in\R^n\), and put
\[
   m:=|G\mathbin{\triangle}B_1(z)|.
\]
If
\[
   m<\min\{c_dr_d^n,\omega_n(r_d/2)^n\},
\tag{4.32}\label{eq:m-small-H}
\]
then
\[
   d_H\bigl(\partial G,\partial B_1(z)\bigr)
   \le C_H(n,c_d)\,m^{1/n},
\tag{4.33}\label{eq:hausdorff-from-measure}
\]
where one may take
\[
   C_H(n,c_d):=\max\{2c_d^{-1/n},\,4\omega_n^{-1/n}\}.
\tag{4.34}\label{eq:CH}
\]
\end{lemma}

\begin{proof}
Translations do not affect the assertion, so take \(z=0\).

\emph{Boundary points of \(G\) are close to the sphere.}
Let \(x\in\partial G\), and set \(d=\operatorname{dist}(x,\partial B_1)\).
If \(d=0\), the desired estimate for \(x\) is immediate, so suppose
\(d>0\).
If \(d>2r_d\), then \(B_{r_d}(x)\) is entirely contained either in \(B_1\)
or in \(\R^n\setminus B_1\).  In the first case the second density estimate
in \eqref{eq:density-input} gives
\[
   m\ge|B_{r_d}(x)\setminus G|\ge c_dr_d^n;
\]
in the second case the first density estimate gives the same conclusion.
Both contradict \eqref{eq:m-small-H}.  Thus \(d\le2r_d\).
The ball \(B_{d/2}(x)\) is again entirely on one side of \(\partial B_1\).
Apply the appropriate density inequality at radius \(d/2\le r_d\):
\[
   m\ge c_d(d/2)^n.
\]
It follows that
\[
   \operatorname{dist}(x,\partial B_1)
   \le2c_d^{-1/n}m^{1/n}.
\tag{4.35}\label{eq:G-to-B}
\]

\emph{Boundary points of the sphere are close to \(G\).}
Let \(y\in\partial B_1\), and set
\(d=\operatorname{dist}(y,\partial G)\).
If \(d=0\), the desired estimate for \(y\) is immediate, so suppose
\(d>0\).
If \(d>2r_d\), the connected ball \(B_{2r_d}(y)\), which is disjoint from
\(\partial G\), lies, up to a null set, entirely in \(G\) or entirely in
its complement.
The two balls
\[
   B_{r_d/2}\bigl((1-r_d/2)y\bigr)\subset B_1\cap B_{2r_d}(y),
\]
\[
   B_{r_d/2}\bigl((1+r_d/2)y\bigr)
   \subset(\R^n\setminus B_1)\cap B_{2r_d}(y)
\]
show, in either alternative, that
\[
   m\ge\omega_n(r_d/2)^n,
\]
contrary to \eqref{eq:m-small-H}.  Hence \(d\le2r_d\).
Now \(B_{d/2}(y)\) is disjoint from \(\partial G\), so it lies, up to a
null set, wholly in \(G\) or wholly in its complement.  The same inner/outer ball construction,
with radius \(d/4\), gives
\[
   m\ge\omega_n(d/4)^n.
\]
Therefore
\[
   \operatorname{dist}(y,\partial G)
   \le4\omega_n^{-1/n}m^{1/n}.
\tag{4.36}\label{eq:B-to-G}
\]
Combining \eqref{eq:G-to-B} and \eqref{eq:B-to-G} proves
\eqref{eq:hausdorff-from-measure}.
\end{proof}

\begin{corollary}[The deterministic Plan~1 Hausdorff step]\label{cor:FJ-Hausdorff}
Assume the selected normalised set \(\widetilde F\) in
Theorem~\ref{thm:FJ-selection} satisfies the density hypotheses of
Lemma~\ref{lem:FJ-density-Hausdorff}.  Let \(z_{\widetilde F}\) be a
Fraenkel centre supplied by the constructive quantitative isoperimetric
inequality:
\[
   |\widetilde F\mathbin{\triangle}B_1(z_{\widetilde F})|^2
   \le C_{\rm iso}(n)\Def(\widetilde F).
\tag{4.37}\label{eq:FMP-selected}
\]
When \(\Def(\widetilde F)\) is below the explicit threshold which makes
\eqref{eq:m-small-H} hold, one has
\[
   d_H\bigl(\partial\widetilde F,\partial B_1(z_{\widetilde F})\bigr)
   \le C_H(n,c_d)\,C_{\rm iso}(n)^{1/(2n)}
       \Def(\widetilde F)^{1/(2n)}.
\tag{4.38}\label{eq:selected-Hausdorff}
\]
\end{corollary}

\begin{proof}
Set
\[
   m=|\widetilde F\mathbin{\triangle}B_1(z_{\widetilde F})|.
\]
Estimate \eqref{eq:FMP-selected} gives
\[
   m^{1/n}\le C_{\rm iso}(n)^{1/(2n)}
                 \Def(\widetilde F)^{1/(2n)}.
\]
Insert this inequality into \eqref{eq:hausdorff-from-measure}.
\end{proof}

\subsection{Conditional closure for the potential-deficit attempt}

The next proposition records which part of Plan~1 becomes easier here.  The
final closure in Section~\ref{sec:graphs} only needs a
\(W^{1,\infty}\)-small radial graph; it does not require Plan~1's later
\(C^{2,\gamma}\) Schauder interpolation step.

\begin{proposition}[Conditional \(\mathcal C\)-selection closure]\label{prop:FJ-conditional-close}
Fix \(R>1\).  Suppose there are computable constants
\[
   q_{\rm ent}=q_{\rm ent}(n,R)>0,\qquad
   d_{\rm ent}=d_{\rm ent}(n,R)>0
\tag{4.39}\label{eq:entry-constants}
\]
with \(d_{\rm ent}\le\Lambda\omega_n/4\), and with the following
graph-entry property:

\smallskip
\emph{Whenever \(E_0\subset B_R\), \(|E_0|=|B_1|\),
\(\mathcal C(E_0)>0\),}
\[
   \Def(E_0)\le d_{\rm ent},\qquad
   Q_{\mathcal C}(E_0)\le q_{\rm ent},
\tag{4.40}\label{eq:entry-input}
\]
\emph{the selected and volume-normalised set \(\widetilde F\) of
Theorem~\ref{thm:FJ-selection}, after translation by its barycentre, has the
form \(E_u\) in \eqref{eq:radial-set}, satisfies \eqref{eq:graph-balanced},
and obeys}
\[
   \norm{u}_{W^{1,\infty}(S)}\le\delta_*(n)=\frac1{100n^3}.
\tag{4.41}\label{eq:entry-output}
\]

Then every \(E\subset B_R\) with \(|E|=|B_1|\) satisfies
\[
   \beta(E)^2\le C_{\rm bdd}(n,R)\Def(E),
\tag{4.42}\label{eq:bounded-effective-FJ}
\]
where one may take
\[
   q_*:=\min\left\{\frac1{144},q_{\rm ent},\frac18\right\},
\qquad
   C_{\rm bdd}(n,R)
   :=1+\max\left\{\frac1{q_*},
                   \frac{\Per(B_1)}{d_{\rm ent}}\right\}.
\tag{4.43}\label{eq:C-bdd}
\]
\end{proposition}

\begin{proof}
If \(\mathcal C(E)=0\), then \(\beta(E)^2=\Def(E)\), so assume
\(\mathcal C(E)>0\).  If \(\Def(E)>d_{\rm ent}\), then
\(\mathcal C(E)\le\Per(B_1)\) by \eqref{eq:DC-general}, and therefore
\[
   \mathcal C(E)
   \le\frac{\Per(B_1)}{d_{\rm ent}}\Def(E).
\tag{4.44}\label{eq:large-D-C}
\]
If \(Q_{\mathcal C}(E)\ge q_*\), then
\[
   \mathcal C(E)\le q_*^{-1}\Def(E).
\tag{4.45}\label{eq:large-q-C}
\]
It remains to exclude simultaneously
\[
   \Def(E)\le d_{\rm ent},\qquad
   Q_{\mathcal C}(E)<q_*.
\tag{4.46}\label{eq:bad-regime}
\]
Apply Theorem~\ref{thm:FJ-selection}; its hypotheses hold because
\(q_*\le1/144\) and \(d_{\rm ent}\le\Lambda\omega_n/4\).  The selected normalized set
\(\widetilde F\) obeys
\[
   Q_{\mathcal C}(\widetilde F)\le2Q_{\mathcal C}(E)<\frac14.
\tag{4.47}\label{eq:q-upper-selected}
\]
On the other hand, graph entry and Corollary~\ref{cor:graph-simple} give
\[
   \Def(\widetilde F)+\mathcal C(\widetilde F)
   =\beta(\widetilde F)^2
   \le3\Def(\widetilde F),
\]
so
\[
   \mathcal C(\widetilde F)\le2\Def(\widetilde F),
   \qquad
   Q_{\mathcal C}(\widetilde F)\ge\frac12,
\tag{4.48}\label{eq:q-lower-selected}
\]
a contradiction.  Hence either \eqref{eq:large-D-C} or
\eqref{eq:large-q-C} holds.  Adding \(\Def(E)\) to the resulting control of
\(\mathcal C(E)\), and using \eqref{eq:beta-DC}, proves
\eqref{eq:bounded-effective-FJ}--\eqref{eq:C-bdd}.
\end{proof}

\subsection{What the potential-deficit attempt leaves unresolved}

Proposition~\ref{prop:FJ-conditional-close} reduces a bounded effective FJ
inequality to the single graph-entry output \eqref{eq:entry-output}.
The selection map and the quotient preservation needed for that reduction
have now been proved in Theorem~\ref{thm:FJ-selection}.  A complete
constructive graph-entry theorem along this first route would proceed as
follows.  The direct-\(\beta\) regularisation below supersedes steps (G1)--(G3)
by proving an additive almost-minimality estimate.

\begin{enumerate}[label=\textnormal{(G\arabic*)},leftmargin=2.8em]
\item Use \eqref{eq:near-min}, with \(\tau=q^{1/4}\) below an explicit
threshold, to prove two-sided density estimates for \(F\) with constants
depending only on \(n,R,\Lambda\).
\item After volume normalisation, note from
\eqref{eq:selected-def-volume} that \(|F|\ge3\omega_n/4\) and hence the
normalising factor satisfies
\[
   (4/5)^{1/n}\le\lambda\le(4/3)^{1/n}.
\]
If \(F\) has the density estimate in (G1) up to radius \(r_d\), then
\(\widetilde F=\lambda F\) has the same density coefficient up to radius
\(\lambda r_d\), by applying the estimate for \(F\) at radius \(r/\lambda\).
Corollary~\ref{cor:FJ-Hausdorff} then upgrades the constructive quantitative
isoperimetric inequality and this scaled density estimate to the Hausdorff
annular trapping estimate \eqref{eq:selected-Hausdorff}.
\item Establish an effective improvement-of-flatness theorem for minimisers
of \eqref{eq:J-FJ}, using the near-minimality constants displayed in
\eqref{eq:near-min}.  The required conclusion is only a global
\(W^{1,\infty}\) radial graph with norm at most \(\delta_*(n)\).
\item Translate to the barycentre and quantify that the radial graph norm
remains below \(\delta_*(n)\).  Volume normalisation is already harmless
because \(Q_{\mathcal C}\) is exactly scale invariant.
\item To globalise this diagnostic \(\mathcal C\)-selection route one would
need a truncation lemma controlling \(Q_{\mathcal C}\).  The direct-\(\beta\)
route below instead uses Lemma~\ref{lem:computable-confinement}.
\end{enumerate}

\begin{remark}[Comparison with the Plan~1 gaps]\label{rem:Plan1-comparison}
Plan~1 required graph entry followed by a small \(C^{2,\gamma}\) bootstrap
because its final Saint--Venant expansion used that norm.  In the present
FJ adaptation, the fully proved closure theorem,
Corollary~\ref{cor:graph-simple}, only assumes small
\(W^{1,\infty}\).  Therefore the Schauder-interpolation gap does not appear.
For this diagnostic \(\mathcal C\)-selection attempt alone, the unresolved
analytic point is the near-minimiser-to-graph entry in (G1)--(G3), since its
error has perimeter order.  The direct-\(\beta\) regularisation below removes
that obstruction and supplies its graph-entry computability certificate.
Its computable confinement step then gives the unrestricted theorem.
\end{remark}

\subsection{Regularising selection by penalising the true excess}
\label{ssec:true-beta-selection}

The preceding attempt selected a set while preserving
\(\mathcal C\).  It exposed the obstruction, but it is not the most effective
regularisation.  The original Fusco--Julin construction \cite{FJ} instead
preserves \(\beta^2\) itself.  At first sight this appears worse, because
\(\beta^2=\Per-(n-1)M\) contains the perimeter and the singular potential
\(M\).  The point is that the argument is performed in two stages:
\[
\begin{split}
   &\text{penalisation of \(\beta^2\)}
   \Longrightarrow \text{area quasi-minimality and density},\\
   &\text{density plus small deficit}
   \Longrightarrow \text{annular trapping and centre separation},\\
   &\text{centre separation}
   \Longrightarrow \text{additive perimeter almost-minimality}.
\end{split}
\tag{4.49}\label{eq:regularising-chain}
\]
The last implication is where the maximiser in \(M\) becomes regular enough:
on all local variations touching the boundary, the kernel
\(|x-y|^{-1}\) is uniformly bounded because every maximising centre \(y\)
lies strictly inside the trapped annulus.

For this subsection set
\[
   b(U):=\beta(U)^2=\Per(U)-(n-1)M(U).
\tag{4.50}\label{eq:bU}
\]
Fix \(R>3\), a container \(B_R\), a number \(\Lambda>n\), and set
\[
   a:=\frac14.
\tag{4.51}\label{eq:a-quarter}
\]
For a bounded input one may translate the configuration without changing
\(b\) or \(\Def\).  Lemma~\ref{lem:computable-confinement} below supplies
bounded inputs from arbitrary finite-perimeter sets.

\begin{lemma}[The volume penalty baseline]\label{lem:volume-baseline}
For every finite-perimeter \(U\subset B_R\), put
\[
   H_\Lambda(U):=
   \Per(U)+\Lambda\bigl||U|-\omega_n\bigr|-\Per(B_1).
\tag{4.52}\label{eq:H-Lambda}
\]
Then
\[
   H_\Lambda(U)
   \ge \Def(U)+(\Lambda-n)\bigl||U|-\omega_n\bigr|
   \ge0.
\tag{4.53}\label{eq:H-lower}
\]
\end{lemma}

\begin{proof}
Write \(m=|U|\) and \(\alpha=(n-1)/n\).  By the isoperimetric inequality,
\[
   \Per(U)\ge p(m)=n\omega_n^{1/n}m^\alpha.
\]
If \(m\ge\omega_n\), then \(p(m)\ge p(\omega_n)=\Per(B_1)\), and
\[
   p(m)+\Lambda(m-\omega_n)-\Per(B_1)
   \ge\Lambda(m-\omega_n)
   \ge(\Lambda-n)(m-\omega_n).
\]
If \(0\le m\le\omega_n\), set \(t=m/\omega_n\).  Since
\(0\le t\le1\) and \(\alpha<1\), \(t^\alpha\ge t\).  Therefore
\[
\begin{aligned}
   p(m)+\Lambda(\omega_n-m)-\Per(B_1)
   &=\omega_n\{n(t^\alpha-1)+\Lambda(1-t)\}\\
   &\ge(\Lambda-n)(\omega_n-m).
\end{aligned}
\]
Subtracting \(p(m)\) from \(\Per(U)\) gives \eqref{eq:H-lower}.
\end{proof}

\begin{theorem}[Regularising single-set FJ selection]
\label{thm:true-beta-selection}
Let \(E_0\subset B_R\) satisfy
\[
   |E_0|=\omega_n,\qquad
   b_0:=b(E_0)>0,\qquad
   d_0:=\Def(E_0),\qquad
   q_0:=\frac{d_0}{b_0}.
\tag{4.54}\label{eq:true-beta-input}
\]
Assume
\[
   q_0\le\frac1{16},
   \qquad
   d_0\le\frac{\Lambda-n}{2}\,\omega_n.
\tag{4.55}\label{eq:true-beta-small}
\]
Then the functional
\[
   \mathcal I_{E_0}(U)
   :=
   \Per(U)+\Lambda\bigl||U|-\omega_n\bigr|
      +a\,|b(U)-b_0|,
   \qquad U\subset B_R,
\tag{4.56}\label{eq:true-beta-functional}
\]
has a minimiser \(F\).  Every minimiser satisfies
\[
   \Def(F)\le d_0,\qquad
   \bigl||F|-\omega_n\bigr|
      \le\frac{d_0}{\Lambda-n},
\tag{4.57}\label{eq:true-selected-energy}
\]
\[
   |b(F)-b_0|\le4d_0,\qquad
   b(F)\ge(1-4q_0)b_0\ge\frac34b_0,
\tag{4.58}\label{eq:true-selected-beta}
\]
and
\[
   \frac{\Def(F)}{b(F)}
      \le\frac{q_0}{1-4q_0}\le\frac43q_0.
\tag{4.59}\label{eq:true-selected-ratio}
\]
If \(s=(|F|/\omega_n)^{1/n}\) and
\(\widetilde F=s^{-1}F\), then
\[
   |\widetilde F|=\omega_n,\qquad
   \frac{\Def(\widetilde F)}{b(\widetilde F)}
      =\frac{\Def(F)}{b(F)},\qquad
   \Def(\widetilde F)\le2d_0.
\tag{4.60}\label{eq:true-selected-normalised}
\]
\end{theorem}

\begin{proof}
For existence, let \(U_j\) be a minimising sequence.  Comparison with
\(E_0\) bounds \(\Per(U_j)\), since all terms in
\eqref{eq:true-beta-functional} except the perimeter are nonnegative.
After passing to a subsequence, \(U_j\to F\) in \(L^1(B_R)\) and
\(\Per(F)\le\liminf_j\Per(U_j)\).  Lemma~\ref{lem:M-cont} implies
\(M(U_j)\to M(F)\).  For fixed \(c\), the function
\[
   t\longmapsto t+a|t-c|
\]
is nondecreasing when \(0<a<1\).  Applying this with
\(c=(n-1)M(U_j)+b_0\), which converges to
\((n-1)M(F)+b_0\), proves lower semicontinuity of
\(\Per(U)+a|b(U)-b_0|\).  The volume term is continuous; hence \(F\)
is a minimiser.

Minimality against \(E_0\), followed by subtraction of \(\Per(B_1)\),
gives
\[
   H_\Lambda(F)+a|b(F)-b_0|\le d_0.
\tag{4.61}\label{eq:true-min-compare}
\]
Lemma~\ref{lem:volume-baseline} yields
\eqref{eq:true-selected-energy}; since \(a=1/4\), it also yields
\(|b(F)-b_0|\le4d_0\).  Dividing the latter estimate by \(b_0\)
proves \eqref{eq:true-selected-beta}, and then
\[
   \frac{\Def(F)}{b(F)}
   \le\frac{d_0}{(1-4q_0)b_0}
   =\frac{q_0}{1-4q_0},
\]
which proves \eqref{eq:true-selected-ratio}.

From \eqref{eq:true-selected-energy} and
\eqref{eq:true-beta-small}, \(|F|\ge\omega_n/2\), hence
\(s\ge2^{-1/n}\).  Both \(\Def\) and \(b\) scale by \(s^{n-1}\);
therefore the quotient in \eqref{eq:true-selected-normalised} is
scale invariant.  Finally,
\[
   \Def(\widetilde F)
   =s^{1-n}\Def(F)
   \le2^{(n-1)/n}d_0\le2d_0.
\]
\end{proof}

\begin{proposition}[First regularisation stage: area quasi-minimality]
\label{prop:area-quasi}
Let \(F\) be a minimiser in Theorem~\ref{thm:true-beta-selection}.  If
\[
   F\mathbin{\triangle}V\Subset B_r(x),
\tag{4.62}\label{eq:beta-local-comp}
\]
then
\[
   \Per(F;B_r(x))
   \le3\Per(V;B_r(x))
        +2\Lambda|F\mathbin{\triangle}V|.
\tag{4.63}\label{eq:coarse-local}
\]
Consequently, for
\[
   r_a:=\min\left\{1,\frac{n}{4\Lambda}\right\},
\tag{4.64}\label{eq:ra}
\]
and \(0<r\le r_a\),
\[
   \Per(F;B_r(x))\le7\Per(V;B_r(x)).
\tag{4.65}\label{eq:area-seven}
\]
\end{proposition}

\begin{proof}
First suppose \(B_r(x)\Subset B_R\), so that \(V\) itself is an
admissible competitor in \eqref{eq:true-beta-functional}.
Minimality and the triangle inequality for the last two terms in
\eqref{eq:true-beta-functional} give
\[
   \Per(F)\le\Per(V)+\Lambda|F\mathbin{\triangle}V|
                  +a|b(F)-b(V)|.
\tag{4.66}\label{eq:min-with-bdiff}
\]
If \(b(F)\ge b(V)\), select a centre \(y_V\) realising \(b(V)\);
if \(b(V)>b(F)\), use a centre \(y_F\) realising \(b(F)\).  In either
case the oscillation integrand
\[
   1-\nu_U(z)\cdot\frac{z-y}{|z-y|}
\]
lies between \(0\) and \(2\).  Since \(F=V\) outside \(B_r(x)\),
\[
   |b(F)-b(V)|
   \le2\{\Per(F;B_r(x))+\Per(V;B_r(x))\}.
\tag{4.67}\label{eq:b-local-coarse}
\]
Insert \(a=1/4\) into \eqref{eq:min-with-bdiff}, cancel the
perimeters outside \(B_r(x)\), and move
\(\frac12\Per(F;B_r(x))\) to the left.  This proves
\eqref{eq:coarse-local}.

Let \(A=F\mathbin{\triangle}V\).  Since \(A\subset B_r(x)\), the
isoperimetric inequality and \(P(A)\le\Per(F;B_r(x))+\Per(V;B_r(x))\)
give
\[
   |A|\le\frac r n
       \{\Per(F;B_r(x))+\Per(V;B_r(x))\}.
\tag{4.68}\label{eq:A-perimeters}
\]
For \(r\le n/(4\Lambda)\), inserting this estimate in
\eqref{eq:coarse-local} gives
\[
   \Per(F;B_r(x))
   \le3\Per(V;B_r(x))
      +\frac12\{\Per(F;B_r(x))+\Per(V;B_r(x))\}.
\]
Rearrangement proves \eqref{eq:area-seven}.

If \(B_r(x)\) is not contained in the container, apply the preceding
argument to \(W=V\cap B_R\), which is an admissible competitor because
\(F\subset B_R\).  Since \(B_R\subset V\cup B_R\), the isoperimetric
inequality gives
\[
   \Per(B_R)\le\Per(V\cup B_R).
\]
The perimeter lattice inequality
\[
   \Per(V\cap B_R)+\Per(V\cup B_R)
      \le\Per(V)+\Per(B_R)
\]
then implies \(\Per(W)\le\Per(V)\).  The same cancellation outside
\(B_r(x)\) shows \(\Per(W;B_r(x))\le\Per(V;B_r(x))\), while
\(|F\mathbin{\triangle}W|\le|F\mathbin{\triangle}V|\).
Thus \eqref{eq:coarse-local} and \eqref{eq:area-seven} also hold for
balls meeting \(\partial B_R\).
\end{proof}

\begin{lemma}[Explicit density from the coarse stage]\label{lem:area-density}
Let \(F\) satisfy \eqref{eq:area-seven} and take the representative whose
boundary is \(\operatorname{spt}|D{\bf1}_F|\).  Then, for every
\(x\in\partial F\) and \(0<r\le r_a\),
\[
   |F\cap B_r(x)|\ge c_a r^n,\qquad
   |B_r(x)\setminus F|\ge c_a r^n,
   \qquad
   c_a:=\frac{\omega_n}{8^n}.
\tag{4.69}\label{eq:area-density}
\]
\end{lemma}

\begin{proof}
Set \(m(r)=|F\cap B_r(x)|\).  For almost every \(r\), comparison with
the set obtained by removing \(F\) from \(B_r(x)\), using an arbitrarily
thin collar to meet the compact-containment condition, gives
\[
   \Per(F;B_r(x))\le7m'(r).
\]
The perimeter of \(F\cap B_r(x)\) is at most
\(\Per(F;B_r(x))+m'(r)\le8m'(r)\).  Thus the isoperimetric
inequality gives
\[
   n\omega_n^{1/n}m(r)^{(n-1)/n}\le8m'(r)
\]
for almost every \(r\).  Since \(x\) belongs to the support of the
perimeter, \(m(r)>0\) for every \(r>0\), and integration of
\[
   \frac d{dr}m(r)^{1/n}
   \ge\frac{\omega_n^{1/n}}8
\]
from \(0\) to \(r\) yields the first inequality in
\eqref{eq:area-density}.  Apply the same argument to the complement of
\(F\), whose local perimeter satisfies the identical comparison
inequality, to obtain the second one.
\end{proof}

\begin{lemma}[Quantitative localisation of potential centres]
\label{lem:potential-centre}
Let
\[
   K_n:=\frac{n}{n-1}\omega_n^{1/n}
\]
be as in Lemma~\ref{lem:M-cont}, and define the computable positive number
\[
\begin{aligned}
   \Gamma_n
   &:=
   \int_{B_1}\frac{dx}{|x|}
     -\int_{B_1}\frac{dx}{|x-\frac14 e_1|}>0,\\
   m_c(n)
   &:=
   \left(\frac{\Gamma_n}{4K_n}\right)^{n/(n-1)}.
\end{aligned}
\tag{4.70}\label{eq:Gamma-mc}
\]
Let \(G\) have finite measure and let \(B_s(z)\), \(s>0\), satisfy
\[
   \frac{|G\mathbin{\triangle}B_s(z)|}{s^n}\le m_c(n).
\tag{4.71}\label{eq:center-mismatch}
\]
Then every maximiser \(y_G\) in the definition of \(M(G)\) satisfies
\[
   |y_G-z|<\frac{s}{4}.
\tag{4.72}\label{eq:center-localised}
\]
\end{lemma}

\begin{proof}
The function \(y\mapsto\int_G|x-y|^{-1}\,dx\) is continuous, tends to
zero as \(|y|\to\infty\), and therefore admits a maximiser.  By
translation and scaling it is enough to take \(s=1\), \(z=0\).
The potential of \(B_1\) is a strictly radially decreasing function of
the centre distance; consequently, if \(|y|\ge1/4\),
\[
   \int_{B_1}\frac{dx}{|x|}
      -\int_{B_1}\frac{dx}{|x-y|}
   \ge\Gamma_n.
\]
Maximality of \(y_G\), and the rearrangement estimate used in
Lemma~\ref{lem:M-cont}, give
\[
\begin{aligned}
   \Gamma_n
   &\le
   \left|\int_{B_1}\frac{dx}{|x|}
          -\int_G\frac{dx}{|x|}\right|
   +\left|\int_G\frac{dx}{|x-y_G|}
          -\int_{B_1}\frac{dx}{|x-y_G|}\right|\\
   &\le 2K_n|G\mathbin{\triangle}B_1|^{(n-1)/n}
   \le\frac{\Gamma_n}{2},
\end{aligned}
\]
a contradiction.  Scaling gives \eqref{eq:center-localised}.
\end{proof}

\begin{theorem}[Second regularisation stage: additive almost-minimality]
\label{thm:additive-almost-min}
Let \(F\) be selected by Theorem~\ref{thm:true-beta-selection}, and let
\(\widetilde F=s^{-1}F\) be its volume normalisation.  Assume that a
constructive quantitative isoperimetric inequality supplies
\[
   |\widetilde F\mathbin{\triangle}B_1(z)|^2
   \le C_{\rm iso}(n)\Def(\widetilde F).
\tag{4.73}\label{eq:FMP-true-selected}
\]
Let \(C_H=C_H(n,c_a)\) be the constant in
Lemma~\ref{lem:FJ-density-Hausdorff}, with density radius
\[
   \widetilde r_a:=\frac23 r_a,
\]
and set
\[
\begin{aligned}
   \overline m
   &:=
   \frac12\min\left\{
        c_a\widetilde r_a^n,\,
        \omega_n(\widetilde r_a/2)^n,\,
        (8C_H)^{-n},\,
        m_c(n)\right\},\\
   d_{\rm ann}
   &:=
   \min\left\{
      \frac{\overline m^{\,2}}{2C_{\rm iso}(n)},\,
      \frac{\Lambda-n}{2}\omega_n
      \right\},\\
   r_{\rm reg}
   &:=
   \min\left\{
      r_a,\frac1{32},
      \frac12\left(\frac{m_c(n)}{2\omega_n}\right)^{1/n}
      \right\}.
\end{aligned}
\tag{4.74}\label{eq:annular-constants}
\]
If \(d_0\le d_{\rm ann}\), then \(F\) is an
\((\Lambda_{\rm reg},r_{\rm reg})\)-almost minimiser of perimeter in
the standard volume-error sense:
for every competitor as in \eqref{eq:beta-local-comp} with
\(0<r\le r_{\rm reg}\),
\[
\begin{aligned}
   \Per(F;B_r(x))
   &\le\Per(V;B_r(x))
       +\Lambda_{\rm reg}|F\mathbin{\triangle}V|\\
   &\le\Per(V;B_r(x))+\Omega_{\rm reg}r^n,\\
   \Lambda_{\rm reg}&:=\frac43\bigl(\Lambda+n-1\bigr),
   \qquad
   \Omega_{\rm reg}:=\omega_n\Lambda_{\rm reg}.
\end{aligned}
\tag{4.75}\label{eq:additive-almost}
\]
The normalised set \(\widetilde F\) satisfies the same estimate with
radius \((2/3)r_{\rm reg}\) and constant at most
\((3/2)\Lambda_{\rm reg}\) in the volume-error form.
\end{theorem}

\begin{proof}
By \eqref{eq:true-selected-normalised} and
\eqref{eq:FMP-true-selected},
\[
   m:=|\widetilde F\mathbin{\triangle}B_1(z)|
   \le(2C_{\rm iso}(n)d_0)^{1/2}
   \le\overline m.
\tag{4.76}\label{eq:m-selected-small}
\]
Proposition~\ref{prop:area-quasi} and Lemma~\ref{lem:area-density}
are invariant under scaling; since
\eqref{eq:true-selected-energy} gives
\[
   2^{-1/n}\le s\le(3/2)^{1/n}\le\frac32,
\]
the normalised set has density coefficient \(c_a\) at least up to
radius \(\widetilde r_a\).  Lemma~\ref{lem:FJ-density-Hausdorff} and
\eqref{eq:annular-constants} yield
\[
   d_H(\partial\widetilde F,\partial B_1(z))\le\frac18.
\tag{4.77}\label{eq:one-eighth-annulus}
\]
After scaling back, every point of \(\partial F\) has distance from
\(z_F:=sz\) between \(7s/8\) and \(9s/8\).  Moreover
\eqref{eq:m-selected-small} and Lemma~\ref{lem:potential-centre}
give
\[
   |y_F-z_F|<\frac{s}{4}
\tag{4.78}\label{eq:yF-inner}
\]
for every potential centre of \(F\).

First take \(V\subset B_R\) as in the statement.  If
\(\Per(F;B_r(x))=0\), \eqref{eq:additive-almost} is immediate.
Otherwise \(B_r(x)\) meets \(\partial F\).  If
\(A=F\mathbin{\triangle}V\), then
\[
   \frac{|V\mathbin{\triangle}B_s(z_F)|}{s^n}
   \le m+\frac{|A|}{s^n}
   \le\frac12m_c(n)+\omega_n(2r)^n
   \le m_c(n).
\]
Lemma~\ref{lem:potential-centre} therefore gives
\(|y_V-z_F|<s/4\) for every potential centre of \(V\).
For \(w\in A\), choose a boundary point of \(F\) in \(B_r(x)\).
Using \(s\ge1/2\), \(r\le1/32\), \eqref{eq:one-eighth-annulus},
and either \(y=y_F\) or \(y=y_V\), we obtain
\[
   |w-y|
   \ge s\left(1-\frac18-\frac14\right)-2r
   \ge\frac14.
\tag{4.79}\label{eq:kernel-separated}
\]
Evaluating \(M(F)\) at a maximising centre of \(F\) and \(M(V)\) at
a maximising centre of \(V\), and using \eqref{eq:kernel-separated},
gives
\[
   |M(F)-M(V)|\le4|F\mathbin{\triangle}V|.
\tag{4.80}\label{eq:M-local-lip}
\]

Let \(X=\Per(F;B_r(x))-\Per(V;B_r(x))\).  If \(X\le0\), there is
nothing to prove.  If \(X>0\), then
\[
   |b(F)-b(V)|
   \le X+4(n-1)|F\mathbin{\triangle}V|.
\]
Insert this in \eqref{eq:min-with-bdiff}, cancel outside perimeters,
and use \(a=1/4\):
\[
   (1-a)X
   \le\{\Lambda+4a(n-1)\}|F\mathbin{\triangle}V|
   \le\{\Lambda+n-1\}\omega_n r^n.
\]
Since \(1-a=3/4\), this proves both inequalities in
\eqref{eq:additive-almost}.  Finally, scaling a competitor
for \(\widetilde F\) by \(s\) multiplies perimeter by \(s^{n-1}\)
and symmetric-difference volume by \(s^n\); thus the volume-error
almost-minimality constant is multiplied by \(s\le3/2\), while its
admissible radius is at least \((2/3)r_{\rm reg}\).

For a competitor \(V\) not contained in \(B_R\), replace it first by
\(V\cap B_R\).  The perimeter lattice argument at the end of
Proposition~\ref{prop:area-quasi} does not increase either the local
perimeter or the symmetric-difference error; hence the same estimate
holds without requiring \(B_r(x)\Subset B_R\).
\end{proof}

\begin{remark}[What has now been regularised]\label{rem:regularised-centre}
Theorem~\ref{thm:additive-almost-min} removes the singular
\(\tau r^{n-1}\) error in \eqref{eq:near-min}.  It does not do so by
pretending that \(M\) is globally Lipschitz in measure.  Rather, the
selected minimiser first becomes geometrically trapped; after that, all
maximising centres of it and of its local competitors lie a definite
distance away from the varying boundary.  Formula
\eqref{eq:M-local-lip} is then the regularised local first-variation
estimate required by perimeter regularity theory.
\end{remark}

\subsection{Reducing graph entry to one local flatness number}

Theorems~\ref{thm:true-beta-selection} and
\ref{thm:additive-almost-min} have removed the Fusco--Julin-specific
singularity.  We now reduce the remaining regularity work to a single local
epsilon-regularity number for perimeter almost-minimisers.  Existence of such
a number follows from the regularity theorem of Tamanini, in the form
presented for example in \cite[Chapters 24--26]{Maggi}.  For the present
purpose, existence alone is not enough: the harmonic-approximation lemma in
that presentation is proved by compactness.  The lemma below replaces that
one use of compactness by a finite spectral truncation and thereby certifies
computability of the required number.

\begin{definition}[Meaning of computable in this section]\label{def:computable}
A dimensional constant is called \emph{computable} if it can be selected by
a finite list of explicit inequalities from rational numbers, \(\omega_k\),
elementary operations, and finitely many Dirichlet eigenvalues and
eigenfunction bounds on Euclidean balls.  These spectral data are
computable: in one dimension they are trigonometric, and in higher
dimensions they are given by spherical harmonics and zeros of Bessel
functions, which admit certified rational bracketing.  Numerical evaluation
of the resulting expression is not required.
\end{definition}

\begin{definition}[Admissible flatness number]\label{def:flat-number}
Fix \(n\ge2\) and \(0<\ell<1/8\).  A number
\(\varepsilon_{\rm fl}(n,\ell)\in(0,1/16)\) is called
\emph{admissible} if the following statement holds.

Let \(G\) be represented so that
\(\partial G=\operatorname{spt}|D{\bf1}_G|\), and suppose that, in every
ball of radius at most \(r_*\),
\[
   \Per(G;B_t(y))
   \le\Per(V;B_t(y))+\Lambda_*|G\mathbin{\triangle}V|
   \quad\text{whenever }G\mathbin{\triangle}V\Subset B_t(y).
\tag{4.81}\label{eq:local-am-input}
\]
Let \(x\in\partial G\), \(e\in S^{n-1}\), and \(0<16\rho\le r_*\).
If
\[
\begin{gathered}
   \Lambda_*\rho\le\varepsilon_{\rm fl}(n,\ell),\\
   \partial G\cap B_{8\rho}(x)
      \subset
      \{y: |(y-x)\cdot e|\le
                 \varepsilon_{\rm fl}(n,\ell)\rho\},\\
   x-4\rho e\in G^{(1)},\qquad x+4\rho e\in G^{(0)},
\end{gathered}
\tag{4.82}\label{eq:local-flat-input}
\]
then \(\partial G\cap B_\rho(x)\) is a single Lipschitz graph in
direction \(e\), and its Lipschitz constant is at most \(\ell\).
Here \(G^{(1)}\) and \(G^{(0)}\) denote the points of density one and
zero of \(G\), respectively.
\end{definition}

\begin{lemma}[Constructive harmonic approximation]\label{lem:effective-harmonic}
Let \(m\ge1\), \(B=B_1^m\), and \(D=B_{1/2}^m\).  For every
\(\tau\in(0,1)\) there is a computable number
\(\sigma_{\rm har}(m,\tau)>0\) with the following property.  If
\(u\in W^{1,2}(B)\) satisfies
\[
   \int_B|\nabla u|^2\le1,\qquad
   \left|\int_B\nabla u\cdot\nabla\varphi\right|
      \le\sigma_{\rm har}(m,\tau)
          \norm{\nabla\varphi}_{L^\infty(B)}
   \quad\text{for every }\varphi\in C_c^\infty(B),
\]
then there is a harmonic function \(v\) on \(D\) such that
\[
   \int_D|\nabla v|^2\le1,\qquad
   \int_D|u-v|^2\le\tau.
\]
\end{lemma}

\begin{proof}
Let \(v\) be the harmonic replacement of \(u\) in \(D\), and put
\(w=u-v\in W^{1,2}_0(D)\).  Then
\[
   \int_D|\nabla v|^2\le\int_D|\nabla u|^2\le1,\qquad
   \int_D|\nabla w|^2\le1.
\]
Let \((\lambda_j,\phi_j)_{j\ge1}\) be an \(L^2(D)\)-orthonormal
Dirichlet eigenbasis for \(-\Delta\) on \(D\), ordered so that
\(\lambda_j\) is nondecreasing, and set
\[
   L_j:=\norm{\nabla\phi_j}_{L^\infty(D)}.
\]
Each \(\phi_j\), extended by zero outside \(D\), is Lipschitz and can
be mollified inside \(B\) without increasing its gradient bound in
the limit.  Thus the almost-harmonicity assumption gives
\[
   \left|\int_D\nabla w\cdot\nabla\phi_j\right|
      =\left|\int_D\nabla u\cdot\nabla\phi_j\right|
      \le\sigma_{\rm har}(m,\tau)L_j.
\]
Writing \(w=\sum_{j\ge1}a_j\phi_j\), this says
\[
   |a_j|\le \sigma_{\rm har}(m,\tau)\frac{L_j}{\lambda_j}.
\]
Choose a finite \(N=N(m,\tau)\) such that
\(\lambda_{N+1}^{-1}\le\tau/2\), put
\[
   A_N:=\sum_{j=1}^{N}\left(\frac{L_j}{\lambda_j}\right)^2,
   \qquad
   \sigma_{\rm har}(m,\tau)
      :=\min\left\{1,\sqrt{\frac{\tau}{2A_N}}\right\},
\]
with the evident convention if \(A_N=0\).  The low modes and high
modes then satisfy
\[
\begin{aligned}
   \sum_{j=1}^{N}a_j^2&\le\frac{\tau}{2},\\
   \sum_{j>N}a_j^2
      &\le\lambda_{N+1}^{-1}
          \sum_{j>N}\lambda_j a_j^2
       \le\frac{\tau}{2}\int_D|\nabla w|^2
       \le\frac{\tau}{2}.
\end{aligned}
\]
Therefore \(\int_D|u-v|^2=\int_D|w|^2\le\tau\).
Definition~\ref{def:computable} makes every choice in this construction
computable.
\end{proof}

\begin{proposition}[Computability certificate for graph entry]
\label{prop:flat-number-computable}
For every \(n\ge2\) and \(0<\ell<1/8\), there exists a computable
admissible flatness number \(\varepsilon_{\rm fl}(n,\ell)\) in the
sense of Definition~\ref{def:flat-number}.
\end{proposition}

\begin{proof}
The local regularity proof for \((\Lambda_*,r_*)\)-almost minimisers is
the Tamanini iteration presented in
\cite[Theorems 24.1, 24.2, 25.1, 25.3, and 26.3]{Maggi}.
After scaling a ball \(B_\rho(x)\) to a unit cylinder, the error enters
only through the dimensionless number \(\Lambda_*\rho\).  The proof has
the following four inputs:
\[
\begin{array}{ll}
\textnormal{(i)} & \text{the slab and phase hypotheses, which orient a
                       cylinder;}\\
\textnormal{(ii)}& \text{the Caccioppoli, or reverse Poincare, inequality
                       converting height flatness to excess;}\\
\textnormal{(iii)}&\text{Lipschitz approximation and harmonic
                       approximation;}\\
\textnormal{(iv)}& \text{tilt improvement followed by geometric iteration.}
\end{array}
\]
The two phase points in \eqref{eq:local-flat-input}, together with the
absence of boundary outside the slab, identify the upper and lower
phases in the working cylinder.  The comparison proof of (ii) then gives,
after a fixed contraction of that cylinder, an inequality of the form
\[
   \operatorname{exc}
      \le C(n)\left(\varepsilon_{\rm fl}(n,\ell)^2
                    +\Lambda_*\rho\right),
\]
where the numerical value of \(C(n)\) is obtained by carrying out that
comparison estimate.  The constants in (i), (ii), and the
Lipschitz-approximation part of (iii) arise from relative isoperimetry,
Fubini/Chebyshev estimates, and explicit comparison competitors.  The
tilt estimate in (iv) consists
of inserting the harmonic decay estimate into (ii), choosing a fixed
contraction factor, and taking minima of finitely many preceding
smallness conditions.  These operations meet
Definition~\ref{def:computable}.

The compactness step in the printed chain is the harmonic-approximation
lemma used in the tilt argument: approximate harmonicity is converted
into closeness to a harmonic function by contradiction and the compact
embedding \(W^{1,2}\Subset L^2\).  Replace that invocation, after one
fixed interior contraction of the cylinder, by
Lemma~\ref{lem:effective-harmonic} with \(m=n-1\).  Losing the fixed
factor \(1/2\) only changes the contraction factor in (iv): harmonic
decay is then applied on the smaller disk and the next cylinder is
chosen inside it.

For completeness, the parameter-selection algorithm is as follows.
Fix a contraction factor \(\theta\in(0,1/16)\).  Starting with the
desired slope \(\ell\), choose the target harmonic error \(\tau\) so
that harmonic decay plus the Caccioppoli error is at most
\(\theta^{2\alpha}\) times the preceding excess, for some fixed
\(\alpha\in(0,1/2)\).  Choose the corresponding
\(\sigma_{\rm har}(n-1,\tau)\) by Lemma~\ref{lem:effective-harmonic}.
Then take \(\varepsilon_{\rm fl}(n,\ell)\) to be the minimum of the
finitely many height, Lipschitz-approximation, Caccioppoli, and tilt
smallness requirements with that \(\sigma_{\rm har}\).  Iteration gives
a geometric series for the changes of tangent plane, whose sum is made
at most \(\ell\) by the initial choice of the minimum.  This proves the
single-graph conclusion in Definition~\ref{def:flat-number}.

Every operation in this backward choice is finite, except for summing a
geometric series with specified ratio.  Hence
\(\varepsilon_{\rm fl}(n,\ell)\) is computable, although we do not
evaluate it.
\end{proof}

\begin{lemma}[Orientation of a trapped unit-volume set]\label{lem:phase-orient}
Set
\[
   h_{\rm or}(n):=\min\left\{\frac1{32},\frac1{4n}\right\}.
\]
Let \(G\) be a finite-perimeter set with \(|G|=\omega_n\),
\(\partial G=\operatorname{spt}|D{\bf1}_G|\), and
\[
   d_H(\partial G,\partial B_1(z))\le h\le h_{\rm or}(n).
\tag{4.83}\label{eq:trapped-H}
\]
Then, up to null sets,
\[
   B_{1-h}(z)\subset G\subset B_{1+h}(z).
\tag{4.84}\label{eq:oriented-annulus}
\]
\end{lemma}

\begin{proof}
The open connected regions \(B_{1-h}(z)\) and
\(\R^n\setminus\overline{B_{1+h}(z)}\) do not meet \(\partial G\).
Since the boundary is the support of the perimeter measure, the
characteristic function of \(G\) is almost everywhere constant on each
of these regions.  The exterior region cannot belong to \(G\), because
\(|G|<\infty\).  If the interior region did not belong to \(G\), then
\[
   |G|\le |B_{1+h}\setminus B_{1-h}|
       =\omega_n\bigl((1+h)^n-(1-h)^n\bigr).
\]
Since \(nh\le1/4\), Bernoulli's inequality and
\((1+h)^n\le e^{nh}\le e^{1/4}<3/2\) give
\[
   (1+h)^n-(1-h)^n
      <\frac32-\frac34<1.
\tag{4.85}\label{eq:orientation-small}
\]
This contradicts \(|G|=\omega_n\), proving
\eqref{eq:oriented-annulus}.
\end{proof}

\begin{lemma}[Recentering a small radial graph]\label{lem:graph-recentre}
Set
\[
   C_{\rm rec}(n):=32\bigl(1+n2^n\bigr).
\tag{4.86}\label{eq:Crec}
\]
Let \(|E|=\omega_n\) and suppose that, about a centre \(z\),
\[
   \partial E=\{z+(1+v(\theta))\theta:\theta\in S^{n-1}\},
   \qquad
   \norm{v}_{W^{1,\infty}}\le\kappa
   \le\frac{1}{256(1+n2^n)}.
\tag{4.87}\label{eq:recentre-input}
\]
If \(c_E\) is the barycentre of \(E\), then \(E\) is a radial graph
about \(c_E\), with graph function \(u\) satisfying
\[
   \norm{u}_{W^{1,\infty}}
      \le C_{\rm rec}(n)\kappa.
\tag{4.88}\label{eq:recentre-output}
\]
\end{lemma}

\begin{proof}
Translating, take \(z=0\), and put \(a=c_E\).  Polar integration and
\(|E|=\omega_n\) give
\[
   a=\frac{1}{(n+1)\omega_n}
      \int_{S^{n-1}}(1+v)^{n+1}\theta\,d\sigma.
\]
Since \(\int_{S^{n-1}}\theta\,d\sigma=0\), the mean value theorem and
\(|S^{n-1}|=n\omega_n\) yield
\[
   |a|
   \le n(1+\kappa)^n\kappa
   \le n2^n\kappa.
\tag{4.89}\label{eq:bary-shift}
\]
Put
\[
   \eta:=\kappa+|a|
      \le(1+n2^n)\kappa\le\frac1{256}.
\]
The map
\[
   \Theta(\theta):=
   \frac{(1+v(\theta))\theta-a}
        {|(1+v(\theta))\theta-a|}
\]
is quantitatively close to the identity.  Indeed, writing
\(p=(1+v)\theta-a\), \(R=|p|\), and letting \(\tau\) be a unit tangent
vector, one has
\[
   |R-1|\le\eta,\qquad
   |\Theta-\theta|\le\frac{2\eta}{1-\eta}\le3\eta,\qquad
   |Dp[\tau]-\tau|\le2\kappa\le2\eta.
\]
Since
\[
   D\Theta[\tau]
      =R^{-1}(I-\Theta\otimes\Theta)Dp[\tau],
\]
and \(|\Theta\otimes\Theta-\theta\otimes\theta|
\le2|\Theta-\theta|\), it follows that
\[
   |D\Theta[\tau]-\tau|
   \le\frac{2\eta}{1-\eta}
       +\frac{\eta}{1-\eta}+6\eta
   \le16\eta\le\frac1{16}.
\]
Thus \(\Theta\) is a local bi-Lipschitz map.  The homotopy
\[
   \Theta_t(\theta):=
   \frac{(1+t v(\theta))\theta-ta}
        {|(1+t v(\theta))\theta-ta|},
   \qquad 0\le t\le1,
\]
is well-defined because its denominator is at least \(1-\eta>0\);
therefore \(\Theta\) has degree one.  A degree-one local
homeomorphism of \(S^{n-1}\) is a homeomorphism.  Thus the translated
boundary is a radial graph about \(a\), whose radius is
\[
   1+u(\Theta(\theta))
      =|(1+v(\theta))\theta-a|.
\]
Moreover, the displayed derivative bound gives
\(\norm{D\Theta^{-1}}\le2\), and differentiation of the radius gives
\[
   |D R[\tau]|
      \le (1+\kappa)|\Theta-\theta|+\kappa
      \le7\eta.
\]
Consequently
\[
   \norm{u}_{L^\infty}
      \le\eta,
   \qquad
   \norm{\nabla_\tau u}_{L^\infty}
      \le14\eta.
\]
Together with \eqref{eq:bary-shift}, these inequalities imply the
larger stated bound \eqref{eq:recentre-output}.
\end{proof}

\begin{theorem}[Effective graph entry from an admissible flatness number]
\label{thm:graph-entry-flat}
Let \(G\) have unit volume, take the representative for which
\(\partial G=\operatorname{spt}|D{\bf1}_G|\), and suppose that it
satisfies \eqref{eq:local-am-input} with parameters
\(\Lambda_*,r_*\).  Put
\[
   \kappa_*:=\frac{\delta_*(n)}{C_{\rm rec}(n)},\qquad
   \ell_*:=\frac{\kappa_*}{64},\qquad
   \varepsilon_*:=\varepsilon_{\rm fl}(n,\ell_*),
\tag{4.90}\label{eq:graph-target-data}
\]
where \(\varepsilon_{\rm fl}(n,\ell_*)\) is admissible, and define
\[
\begin{aligned}
   \rho_*&:=
     \min\left\{\frac{r_*}{16},\,
       \frac{\varepsilon_*}{512},\,
       \frac{\varepsilon_*}{16(1+\Lambda_*)},\,
       \frac{\kappa_*}{2048},\,\frac1{64}\right\},\\
   h_*&:=
     \min\left\{\frac{\varepsilon_*\rho_*}{16},\,
       \frac{\rho_*}{16},\,
       h_{\rm or}(n)\right\}.
\end{aligned}
\tag{4.91}\label{eq:graph-entry-constants}
\]
If
\[
   d_H(\partial G,\partial B_1(z))\le h_*,
\tag{4.92}\label{eq:graph-entry-H}
\]
then, after translation by the barycentre of \(G\),
\[
   \partial G
    =\{(1+u(\theta))\theta:\theta\in S^{n-1}\},
   \qquad
   \norm{u}_{W^{1,\infty}}\le\delta_*(n).
\tag{4.93}\label{eq:graph-entry-result}
\]
\end{theorem}

\begin{proof}
After translating, take \(z=0\).  The last restriction in
\eqref{eq:graph-entry-constants} permits direct application of
Lemma~\ref{lem:phase-orient}, which fixes the inner and outer phases.

Fix \(x\in\partial G\), choose \(y\in\partial B_1\) with
\(|x-y|\le h_*\), and put \(e=y\).  If
\(p\in\partial G\cap B_{8\rho_*}(x)\), choose
\(\bar p\in\partial B_1\) with \(|p-\bar p|\le h_*\).  Then
\(|\bar p-y|\le8\rho_*+2h_*\le9\rho_*\), and the elementary sphere
identity
\[
   |(\bar p-y)\cdot y|=\frac12|\bar p-y|^2
\]
gives
\[
   |(p-x)\cdot e|
   \le2h_*+\frac{81}{2}\rho_*^2
   \le\varepsilon_*\rho_*.
\tag{4.94}\label{eq:sphere-to-slab}
\]
Moreover, \(x-4\rho_*e\in B_{1-h_*}\) and
\(x+4\rho_*e\notin B_{1+h_*}\), because \(2h_*<4\rho_*\).
Thus the phase assumptions in \eqref{eq:local-flat-input} hold.  Finally,
\(\Lambda_*\rho_*\le\varepsilon_*\) by
\eqref{eq:graph-entry-constants}.  Definition~\ref{def:flat-number}
implies that \(\partial G\cap B_{\rho_*}(x)\) is a single graph of
slope at most \(\ell_*\), for every \(x\in\partial G\).

It remains to pass from local graphs to one radial graph.  The annular
trapping makes the graph direction uniformly transverse to radial rays.
Indeed, if \(q\in\partial G\cap B_{\rho_*}(x)\), then
\[
   \left|\frac{q}{|q|}-e\right|
      \le4(\rho_*+h_*)
      \le\frac{17}{256}<\frac14.
\]
At almost every point of the Lipschitz graph through \(x\), choose its
unit normal \(\nu\) with \(\nu\cdot e>0\).  The slope bound gives
\[
   |\nu-e|\le2\ell_*,
   \qquad
   \left|\nu-\frac{q}{|q|}\right|
      \le4(\ell_*+\rho_*+h_*)<\frac14.
\]
Consequently radial projection from each graph patch to \(S^{n-1}\)
is locally one-to-one.  To see this without an implicit smoothness
argument, suppose that two points of one patch have the form
\(q_i=t_i\theta\), \(t_1\ne t_2\).  Since
\(|\theta-e|<1/4\), one has \(\theta\cdot e>3/4\).  On the other hand,
the graph slope estimate, applied to \(q_2-q_1=(t_2-t_1)\theta\),
would give
\[
   |\theta\cdot e|
      \le \ell_*|\theta-(\theta\cdot e)e|
      \le\ell_*<\frac18,
\]
a contradiction.  In graph coordinates radial projection is
continuous and locally injective; invariance of domain makes it open.
Its image on \(\partial G\) is nonempty and compact, hence also closed;
connectedness of \(S^{n-1}\) shows that it is onto.  Finally, two
global points on the same ray would both lie within \(h_*\) of the
same spherical point and hence at distance less than
\(2h_*<\rho_*\), putting them in one graph patch and yielding the same
contradiction.  Thus
\[
   \partial G=\{(1+v(\theta))\theta:\theta\in S^{n-1}\}.
\]
Here \(\norm{v}_{L^\infty}\le h_*\).  For a radial graph its normal
satisfies
\[
   \frac{|\nabla_\tau v|}{1+v}
      =\frac{|\nu-(\nu\cdot\theta)\theta|}
             {\nu\cdot\theta},
   \qquad \theta=\frac{q}{|q|}.
\]
The preceding estimate, \(1+v\le1+h_*<2\), and
\(\nu\cdot\theta>1/2\) imply
\[
   \norm{\nabla_\tau v}_{L^\infty}
      \le16(\ell_*+\rho_*+h_*)
      <\frac{\kappa_*}{3}.
\]
Furthermore,
\[
   \norm{v}_{L^\infty}
      \le h_*\le\frac{\rho_*}{16}
      \le\frac{\kappa_*}{32768}.
\]
In particular, whether the \(W^{1,\infty}\)-norm is taken as a maximum
or as a sum of these two quantities,
\[
   \norm{v}_{W^{1,\infty}}\le\kappa_*.
\tag{4.95}\label{eq:pre-recentre-graph}
\]
Since \(\delta_*(n)\le1/8\), the definition of \(\kappa_*\) yields
\(\kappa_*\le1/[256(1+n2^n)]\).  Apply
Lemma~\ref{lem:graph-recentre} and
\eqref{eq:graph-target-data} to obtain
\eqref{eq:graph-entry-result}.
\end{proof}

\begin{remark}[What the certificate settles]\label{rem:exact-flat-remaining}
The proof of Theorem~\ref{thm:graph-entry-flat} contains no compactness
step and no unquantified graph patching.  After
Theorem~\ref{thm:additive-almost-min}, its only non-displayed number is
\(\varepsilon_{\rm fl}(n,\ell_*)\), the local height-to-Lipschitz
threshold in Definition~\ref{def:flat-number}.
Proposition~\ref{prop:flat-number-computable} certifies that this number is
computable, without asking that it be numerically evaluated.  Thus no
non-computable regularity constant remains in the bounded
regularised-selection argument.
\end{remark}

\begin{proposition}[Bounded closure through regularised selection]
\label{prop:true-beta-bounded-close}
Let \(\varepsilon_{\rm fl}(n,\ell_*)\) be the computable admissible
number supplied by Proposition~\ref{prop:flat-number-computable}.  In
Theorem~\ref{thm:graph-entry-flat} take
\[
   \Lambda_*:=\frac32\Lambda_{\rm reg},\qquad
   r_*:=\frac23r_{\rm reg},
\]
and let \(h_*\) be given by \eqref{eq:graph-entry-constants}.  Define
\[
   d_{\rm ent}:=\min\left\{d_{\rm ann},\,
      \frac{1}{2C_{\rm iso}(n)}
      \left(\frac{h_*}{C_H(n,c_a)}\right)^{2n}\right\}.
\tag{4.96}\label{eq:d-entry}
\]
Then, in the bounded class used above,
\[
   \beta(E)^2\le C_{\rm reg}(n,R,\Lambda)\Def(E),
\tag{4.97}\label{eq:regularised-bounded-FJ}
\]
where one may take
\[
   C_{\rm reg}
   :=
   \max\left\{16,\,
       2\left(1+\frac{\Per(B_1)}{d_{\rm ent}}\right)\right\}.
\tag{4.98}\label{eq:Creg}
\]
\end{proposition}

\begin{proof}
Let \(d=\Def(E)\) and \(b=\beta(E)^2\).  If \(d>d_{\rm ann}\),
Lemma~\ref{lem:large} gives
\[
   b\le2\{\Per(B_1)+d\}
     \le2\left(1+\frac{\Per(B_1)}{d_{\rm ent}}\right)d.
\]
If \(d_{\rm ent}<d\le d_{\rm ann}\), the same conclusion follows with
\(d_{\rm ent}\) in place of \(d_{\rm ann}\).
If \(d/b\ge1/16\), then \(b\le16d\).  It remains to exclude
\[
   d\le d_{\rm ent},\qquad \frac d b<\frac1{16}.
\]
Apply Theorem~\ref{thm:true-beta-selection} and
Theorem~\ref{thm:additive-almost-min}.  Equations
\eqref{eq:hausdorff-from-measure} and \eqref{eq:d-entry} place
\(\widetilde F\), represented by
\(\partial\widetilde F=\operatorname{spt}|D{\bf1}_{\widetilde F}|\),
within the Hausdorff threshold \(h_*\).  The perimeter comparisons and
the normal oscillation are unchanged by this representative choice, and
Theorem~\ref{thm:graph-entry-flat} places it in the radial-graph regime of
Corollary~\ref{cor:graph-simple}; therefore
\[
   b(\widetilde F)\le3\Def(\widetilde F),
   \qquad
   \frac{\Def(\widetilde F)}{b(\widetilde F)}\ge\frac13.
\]
On the other hand, \eqref{eq:true-selected-ratio} gives
\[
   \frac{\Def(\widetilde F)}{b(\widetilde F)}
   \le\frac{d/b}{1-4d/b}<\frac1{12},
\]
a contradiction.
\end{proof}

\subsection{Computable confinement and the global conclusion}

\begin{lemma}[Computable Fusco--Julin confinement]\label{lem:computable-confinement}
There are computable constants
\[
   R_{\rm red}(n)>3,\qquad
   A_{\rm red}(n)>0,\qquad B_{\rm red}(n)>0
\tag{4.99}\label{eq:reduction-data}
\]
such that, for every finite-perimeter set \(E\subset\R^n\) with
\(|E|=\omega_n\), there is a finite-perimeter set
\(E_0\subset B_{R_{\rm red}(n)}\), after a translation, satisfying
\[
   |E_0|=\omega_n,\qquad
   b(E)\le b(E_0)+A_{\rm red}(n)\Def(E),\qquad
   \Def(E_0)\le B_{\rm red}(n)\Def(E).
\tag{4.100}\label{eq:computable-confinement}
\]
\end{lemma}

\begin{proof}
The geometric construction is Lemma~3.2 of Fusco--Julin
\cite[Lemma~3.2]{FJ}, which in turn repeats the coordinate-slicing
construction of \cite[Lemma~5.1]{FigalliMaggiPratelli}.  We explain why
the constants in that construction are computable.

Fix first a computable dimensional smallness level \(d_{\rm cut}>0\).
If \(d:=\Def(E)\ge d_{\rm cut}\), take \(E_0=B_1\).  Since
\(b(E)\le\Per(E)=\Per(B_1)+d\), one obtains
\[
   b(E)\le
      \left(1+\frac{\Per(B_1)}{d_{\rm cut}}\right)d,
   \qquad \Def(E_0)=0.
\tag{4.101}\label{eq:cut-large}
\]

Suppose now \(d<d_{\rm cut}\).  The slicing construction performs one
cut in each coordinate direction.  At one step, the coarea formula
integrates the slice perimeters over a fixed-dimensional interval; an
averaging choice of two cutting hyperplanes retains the central part,
does not increase perimeter, and removes at most a dimensional multiple
of \(d\) in volume.  Repeating this operation exactly \(n\) times gives
a set \(G\), contained after translation in a rectangular box of
dimensional diameter \(L_{\rm red}(n)\), for which
\[
   G\subset E,\qquad
   \Per(G)\le\Per(E),\qquad
   \omega_n-|G|\le C_{\rm cut}(n)d.
\tag{4.102}\label{eq:cut-output}
\]
Here the cutting width, \(L_{\rm red}(n)\), and \(C_{\rm cut}(n)\) are
obtained by finitely many applications of the displayed slicing
inequalities in the cited proofs.  Thus they are computable.  Decrease
\(d_{\rm cut}\), if necessary, so that
\(C_{\rm cut}(n)d_{\rm cut}\le\omega_n/2\).

Because \(G\subset E\), evaluation of \(M(E)\) at a maximising centre
for \(G\) gives \(M(E)\ge M(G)\).  Therefore
\[
\begin{aligned}
   b(E)-b(G)
   &\le\Per(E)-\Per(G)\\
   &\le d+\Per(B_1)-\Per(G)\\
   &\le d+\Per(B_1)
       \left[1-\left(\frac{|G|}{\omega_n}\right)^{(n-1)/n}\right]\\
   &\le C_1(n)d.
\end{aligned}
\tag{4.103}\label{eq:beta-cut-control}
\]
The third line is the isoperimetric inequality applied to \(G\); the
last line follows from \eqref{eq:cut-output} and the mean value theorem
on \([1/2,1]\).  Hence \(C_1(n)\) is computable.

Set \(\lambda=(\omega_n/|G|)^{1/n}\) and
\(E_0=\lambda G\), followed by translation of its containing box into
a ball.  Then \(1\le\lambda\le2^{1/n}\), so one may take
\[
   R_{\rm red}(n):=\max\{4,2^{1/n}L_{\rm red}(n)\}.
\]
Moreover, since \(b\) scales by \(\lambda^{n-1}\) and is nonnegative,
\[
   b(E)\le b(G)+C_1(n)d\le b(E_0)+C_1(n)d.
\]
Finally,
\[
\begin{aligned}
   \Def(E_0)
   &=\lambda^{n-1}\Per(G)-\Per(B_1)\\
   &\le\{\lambda^{n-1}-1\}\Per(B_1)+\lambda^{n-1}d
    \le C_2(n)d,
\end{aligned}
\tag{4.104}\label{eq:def-cut-control}
\]
again by \eqref{eq:cut-output} and the mean value theorem on
\([1/2,1]\).  Combining the small- and large-deficit branches gives
\eqref{eq:computable-confinement}, with computable
\(A_{\rm red}\) and \(B_{\rm red}\).
\end{proof}

\begin{theorem}[A truly quantitative Fusco--Julin inequality]
\label{thm:global-computable-FJ}
For every \(n\ge2\), there is a computable dimensional constant
\(C_{\rm FJ}^{\rm comp}(n)>0\) such that every finite-perimeter
\(E\subset\R^n\) with \(|E|=|B_\rho|\) satisfies
\[
   \beta(E)^2\le C_{\rm FJ}^{\rm comp}(n)\Def(E).
\tag{4.105}\label{eq:global-computable-FJ}
\]
\end{theorem}

\begin{proof}
By scaling it suffices to consider \(|E|=\omega_n\).  Apply
Lemma~\ref{lem:computable-confinement} and take
\[
   R=R_{\rm red}(n),\qquad \Lambda=2n
\]
in Proposition~\ref{prop:true-beta-bounded-close}.  Since \(E_0\)
belongs to the bounded class of that proposition,
\[
\begin{aligned}
   b(E)
   &\le b(E_0)+A_{\rm red}(n)\Def(E)\\
   &\le C_{\rm reg}\bigl(n,R_{\rm red}(n),2n\bigr)\Def(E_0)
          +A_{\rm red}(n)\Def(E)\\
   &\le
     \left\{A_{\rm red}(n)
       +B_{\rm red}(n)C_{\rm reg}
             \bigl(n,R_{\rm red}(n),2n\bigr)\right\}\Def(E).
\end{aligned}
\]
Thus one may set
\[
   C_{\rm FJ}^{\rm comp}(n)
   :=
   A_{\rm red}(n)
       +B_{\rm red}(n)C_{\rm reg}
             \bigl(n,R_{\rm red}(n),2n\bigr).
\tag{4.106}\label{eq:CFJ-computable}
\]
Every constant on the right is computable by
Lemma~\ref{lem:computable-confinement},
Proposition~\ref{prop:flat-number-computable}, and the displayed
formulae preceding Proposition~\ref{prop:true-beta-bounded-close}.
The scaling laws in Lemma~\ref{lem:scaling} give the assertion for
general \(\rho>0\).
\end{proof}

\begin{corollary}[Original Fusco--Julin formulation]
\label{cor:original-FJ-computable}
For \(|E|=|B_\rho|\), define the scale-invariant same-centre index
\[
   \mathcal A_{\rm FJ}(E)
   :=
   \inf_{z\in\R^n}
   \left\{
      \frac{|E\triangle B_\rho(z)|}{\rho^n}
      +
      \left(
         \frac1{\rho^{n-1}}
         \int_{\partial^*E}
         \left|\nu_E-\frac{x-z}{|x-z|}\right|^2\,d\Hn
      \right)^{1/2}
   \right\}.
\tag{4.107}\label{eq:A-FJ-combined}
\]
There is a computable dimensional constant
\(\widehat C_{\rm FJ}^{\rm comp}(n)>0\) such that every
finite-perimeter set \(E\subset\R^n\) with \(|E|=|B_\rho|\) satisfies
\[
   \mathcal A_{\rm FJ}(E)^2
      \le \widehat C_{\rm FJ}^{\rm comp}(n)\,
          \rho^{1-n}\Def(E).
\tag{4.108}\label{eq:original-FJ-computable}
\]
\end{corollary}

\begin{proof}
This is the normalization of Theorem~1.1 in \cite{FJ}; in particular,
their normal index is
\[
   \beta_{\rm FJ}(E)^2=\rho^{1-n}\beta(E)^2,
\]
with \(\beta\) as in Definition~\ref{def:beta}.  Proposition~1.2 of
\cite{FJ} proves
\[
   \mathcal A_{\rm FJ}(E)
      +\bigl(\rho^{1-n}\Def(E)\bigr)^{1/2}
      \le C_{\rm cmp}(n)\,\beta_{\rm FJ}(E).
\tag{4.109}\label{eq:FJ-comparison}
\]
The constant \(C_{\rm cmp}(n)\) is computable: the proof in \cite{FJ}
uses the radial Gauss--Green identity, rearrangement of
\(|x-z|^{-1}\), and the concavity estimate for the explicit function
\(t\mapsto(1+t)^{(n-1)/n}\) on \([-1,1]\).  These are finite
inequalities with dimensional coefficients and no compactness argument.
Squaring \eqref{eq:FJ-comparison} and applying
Theorem~\ref{thm:global-computable-FJ} gives
\eqref{eq:original-FJ-computable}, for example with
\[
   \widehat C_{\rm FJ}^{\rm comp}(n)
       := C_{\rm cmp}(n)^2 C_{\rm FJ}^{\rm comp}(n).
\]
\end{proof}

\section{Attempt III: torsion superlevels and weighted Reilly identities}
\label{sec:reilly}

For the level-set route to quantitative Faber--Krahn, it is tempting to use
the additional PDE structure: a superlevel of a torsion function carries its
own torsion function.  This section derives the resulting identities from
first principles and records the exact point at which the method ceases to be
unconditional.

\subsection{Set-up and the boundary defects}

Let \(\Om'\subset\R^n\) be a bounded smooth domain and let
\[
   -\Delta v=1\quad\text{in }\Om',\qquad
   v=0\quad\text{on }\partial\Om',\qquad
   v>0\quad\text{in }\Om'.
\tag{5.1}\label{eq:torsion}
\]
Set
\[
   m:=|\Om'|,\qquad P:=\Per(\Om'),\qquad
   T:=\int_{\Om'}v\,dx,
\]
and on \(\partial\Om'\) write
\[
   g:=|\nabla v|.
\]
By the Hopf boundary lemma \(g>0\) on the smooth boundary.  Define
\[
   \alpha:=\int_{\partial\Om'}\frac1g\,d\Hn,\qquad
   \gamma:=\int_{\partial\Om'}g\,d\Hn,
\]
\[
   D_H:=\alpha\gamma-P^2.
\tag{5.2}\label{eq:DH}
\]
If \(|\Om'|=|B_\rho|\), also define
\[
   D_I:=P^2-\Per(B_\rho)^2.
\tag{5.3}\label{eq:DI}
\]

\begin{lemma}[Flux and weighted variance]\label{lem:flux-variance}
For the torsion function in \eqref{eq:torsion},
\[
   \gamma=m.
\tag{5.4}\label{eq:flux}
\]
If \(\bar g:=m/P\), then
\[
   D_H=\frac{P^2}{m}
      \int_{\partial\Om'}\frac{(g-\bar g)^2}{g}\,d\Hn.
\tag{5.5}\label{eq:DH-variance}
\]
\end{lemma}

\begin{proof}
The outer normal derivative satisfies \(\partial_\nu v=-g\).  Integrating
\(-\Delta v=1\) and using the divergence theorem yields
\[
   m=\int_{\Om'}1\,dx
   =-\int_{\Om'}\Delta v\,dx
   =-\int_{\partial\Om'}\partial_\nu v\,d\Hn
   =\int_{\partial\Om'}g\,d\Hn=\gamma.
\]
For the second identity, expand the integrand:
\[
\begin{aligned}
   \int_{\partial\Om'}\frac{(g-\bar g)^2}{g}\,d\Hn
   &=\int_{\partial\Om'}
       \left(g-2\bar g+\frac{\bar g^2}{g}\right)d\Hn \\
   &=m-2\frac mP P+\frac{m^2}{P^2}\alpha \\
   &=-m+\frac{m^2}{P^2}\alpha \\
   &=\frac m{P^2}(\alpha m-P^2)
    =\frac m{P^2}D_H.
\end{aligned}
\]
Multiplication by \(P^2/m\) gives \eqref{eq:DH-variance}.
\end{proof}

\subsection{Three exact identities}

\begin{theorem}[Weighted Reilly--Weinberger and Rellich identities]
\label{thm:reilly-identities}
Let \(v\) be as in \eqref{eq:torsion}.  Define
\[
   Q:=D^2v+\frac1nI.
\]
Then
\[
   2\int_{\Om'}v|Q|^2\,dx
   =\int_{\partial\Om'}g^3\,d\Hn
     -\left(1+\frac2n\right)T.
\tag{5.6}\label{eq:weighted-Reilly}
\]
For every \(z\in\R^n\),
\[
   \int_{\partial\Om'}g^2(x-z)\cdot\nu\,d\Hn=(n+2)T.
\tag{5.7}\label{eq:Rellich}
\]
Combining the preceding identities gives
\[
   2\int_{\Om'}v|Q|^2\,dx
   =
   \int_{\partial\Om'}g^2
      \left[g-\frac1n(x-z)\cdot\nu\right]\,d\Hn.
\tag{5.8}\label{eq:combined}
\]
\end{theorem}

\begin{proof}
Since \(\Delta v=-1\),
\[
\begin{aligned}
   |Q|^2
   &=\left|D^2v+\frac1nI\right|^2 \\
   &=|D^2v|^2+\frac2n\Delta v+\frac1n \\
   &=|D^2v|^2-\frac1n.
\end{aligned}
\tag{5.9}\label{eq:Q-square}
\]
Introduce the Weinberger function
\[
   \Pi:=|\nabla v|^2+\frac2n v.
\]
Differentiating and using that \(\Delta v=-1\) is constant,
\[
\begin{aligned}
   \Delta\Pi
   &=2|D^2v|^2+2\nabla v\cdot\nabla(\Delta v)+\frac2n\Delta v \\
   &=2|D^2v|^2-\frac2n
    =2|Q|^2.
\end{aligned}
\tag{5.10}\label{eq:Pi-lap}
\]
Apply Green's second identity to \(v\) and \(\Pi\):
\[
   \int_{\Om'}(v\Delta\Pi-\Pi\Delta v)\,dx
   =
   \int_{\partial\Om'}
      \bigl(v\partial_\nu\Pi-\Pi\partial_\nu v\bigr)\,d\Hn.
\tag{5.11}\label{eq:Green-two}
\]
On \(\partial\Om'\), \(v=0\), \(\Pi=g^2\), and
\(\partial_\nu v=-g\).  Therefore the right-hand side is
\[
   \int_{\partial\Om'}g^3\,d\Hn.
\tag{5.12}\label{eq:Green-boundary}
\]
On the left-hand side, \eqref{eq:Pi-lap} and \(\Delta v=-1\) give
\[
   2\int_{\Om'}v|Q|^2\,dx+\int_{\Om'}\Pi\,dx.
\]
An integration by parts in the torsion equation yields
\[
   \int_{\Om'}|\nabla v|^2\,dx
   =\int_{\Om'}v\,dx=T,
\]
so
\[
   \int_{\Om'}\Pi\,dx
   =\left(1+\frac2n\right)T.
\]
Together with \eqref{eq:Green-boundary}, this proves
\eqref{eq:weighted-Reilly}.

To prove \eqref{eq:Rellich}, let \(X(x)=x-z\) and introduce the vector field
\[
   Y:=(X\cdot\nabla v)\nabla v-\frac12|\nabla v|^2X.
\]
A coordinate differentiation gives
\[
   \operatorname{div}Y
   =(X\cdot\nabla v)\Delta v+
      \left(1-\frac n2\right)|\nabla v|^2.
\tag{5.13}\label{eq:Rellich-div}
\]
On the boundary, \(\nabla v=-g\nu\), and hence
\[
\begin{aligned}
   Y\cdot\nu
   &=(X\cdot\nabla v)(\nabla v\cdot\nu)
     -\frac12|\nabla v|^2X\cdot\nu \\
   &=(-gX\cdot\nu)(-g)-\frac12g^2X\cdot\nu \\
   &=\frac12g^2X\cdot\nu.
\end{aligned}
\tag{5.14}\label{eq:Rellich-boundary}
\]
Integrating \eqref{eq:Rellich-div}, and using \(\Delta v=-1\), gives
\[
\begin{aligned}
   \int_{\Om'}\operatorname{div}Y\,dx
   &=-\int_{\Om'}X\cdot\nabla v\,dx+
      \left(1-\frac n2\right)T.
\end{aligned}
\]
Since \(v=0\) on the boundary,
\[
   \int_{\Om'}X\cdot\nabla v\,dx
   =\int_{\Om'}\operatorname{div}(vX)\,dx-n\int_{\Om'}v\,dx
   =-nT.
\]
It follows that
\[
   \int_{\Om'}\operatorname{div}Y\,dx
   =\left(n+1-\frac n2\right)T
   =\frac{n+2}{2}T.
\tag{5.15}\label{eq:Rellich-interior}
\]
Equating \eqref{eq:Rellich-boundary} and
\eqref{eq:Rellich-interior}, and multiplying by \(2\), proves
\eqref{eq:Rellich}.
Finally, substitute
\[
   \left(1+\frac2n\right)T
   =\frac1n\int_{\partial\Om'}g^2(x-z)\cdot\nu\,d\Hn
\]
from \eqref{eq:Rellich} into \eqref{eq:weighted-Reilly}; this gives
\eqref{eq:combined}.
\end{proof}

\begin{remark}[Ball check]\label{rem:ball-check}
For \(\Om'=B_\rho(z)\),
\[
   v(x)=\frac{\rho^2-|x-z|^2}{2n},\qquad
   g=\frac{\rho}{n},\qquad Q=0.
\]
Moreover,
\[
   T(B_\rho)=\frac{\omega_n\rho^{n+2}}{n(n+2)}
\]
and
\[
   \int_{\partial B_\rho}g^3\,d\Hn
   =\left(\frac{\rho}{n}\right)^3
     n\omega_n\rho^{n-1}
   =\frac{\omega_n\rho^{n+2}}{n^2}
   =\left(1+\frac2n\right)T(B_\rho).
\]
Thus both sides of \eqref{eq:weighted-Reilly} vanish for a ball, as they
must.
\end{remark}

\subsection{What the boundary defect controls, conditionally}

\begin{lemma}[Third moment under an upper gradient bound]\label{lem:third-moment}
Let the notation be as in \eqref{eq:torsion}--\eqref{eq:DH}.  Assume in
addition that
\[
   0<g\le M_0\qquad\text{on }\partial\Om'.
\tag{5.16}\label{eq:M0}
\]
Then, with \(\bar g=m/P\),
\[
   0\le
   \int_{\partial\Om'}g^3\,d\Hn-\frac{m^3}{P^2}
   \le
   4M_0^2\frac{m}{P^2}D_H.
\tag{5.17}\label{eq:V3-bound}
\]
\end{lemma}

\begin{proof}
The lower bound follows from Jensen's inequality for the convex function
\(t\mapsto t^3\):
\[
   \frac1P\int_{\partial\Om'}g^3\,d\Hn
   \ge\left(\frac1P\int_{\partial\Om'}g\,d\Hn\right)^3
   =\left(\frac mP\right)^3.
\]
For the upper bound, write \(g=\bar g+h\).  Since
\(\int_{\partial\Om'}h\,d\Hn=0\),
\[
\begin{aligned}
   \int_{\partial\Om'}g^3\,d\Hn-\frac{m^3}{P^2}
   &=3\bar g\int_{\partial\Om'}h^2\,d\Hn
     +\int_{\partial\Om'}h^3\,d\Hn.
\end{aligned}
\tag{5.18}\label{eq:V3expand}
\]
Both \(g\) and its average \(\bar g\) lie in \((0,M_0]\), hence
\(|h|\le M_0\).  Therefore
\[
\begin{aligned}
   \int_{\partial\Om'}g^3\,d\Hn-\frac{m^3}{P^2}
   &\le(3M_0+M_0)\int_{\partial\Om'}h^2\,d\Hn \\
   &=4M_0\int_{\partial\Om'}h^2\,d\Hn.
\end{aligned}
\tag{5.19}\label{eq:V3var}
\]
Since \(g\le M_0\),
\[
   \int_{\partial\Om'}h^2\,d\Hn
   \le
   M_0\int_{\partial\Om'}\frac{h^2}{g}\,d\Hn.
\]
Apply \eqref{eq:DH-variance} to obtain
\[
   \int_{\partial\Om'}h^2\,d\Hn
   \le M_0\frac m{P^2}D_H.
\]
Substitution into \eqref{eq:V3var} proves \eqref{eq:V3-bound}.
\end{proof}

Set
\[
   \Dev_w(\Om'):=\int_{\Om'}v\left|D^2v+\frac1nI\right|^2\,dx
\tag{5.20}\label{eq:Devw}
\]
and
\[
   W(\Om'):=
   \left(1+\frac2n\right)T-\frac{m^3}{P^2}.
\tag{5.21}\label{eq:W}
\]
The identity \eqref{eq:weighted-Reilly} can be rewritten exactly as
\[
   2\Dev_w(\Om')
   =
   \left(\int_{\partial\Om'}g^3\,d\Hn-\frac{m^3}{P^2}\right)
   -W(\Om').
\tag{5.22}\label{eq:Dev-decomp}
\]
Lemma~\ref{lem:third-moment} thus controls the first parenthesis by \(D_H\)
provided \eqref{eq:M0} is known.  It does not control the second term.

\subsection{Why this is not yet a proof of an effective FJ estimate}

There are three logically separate gaps.

\begin{enumerate}[label=\textnormal{(R\arabic*)},leftmargin=2.8em]
\item \emph{High-gradient tail.}  The hypothesis \(g\le M_0\) in
Lemma~\ref{lem:third-moment} does not follow from small \(D_H+D_I\) in the
present general level-set setting.  Algebraically, this is unavoidable in
the displayed estimate: \(D_H\) weights \((g-\bar g)^2\) by \(1/g\), whereas
the Reilly identity contains the third moment \(\int g^3\).
\item \emph{The torsional term.}  To use \eqref{eq:Dev-decomp}, one needs an
upper bound for \(-W(\Om')\).  Even if such a bound is obtained from a
torsion deficit, that deficit measures the geometry of all interior
superlevels, not just the boundary level at which \(D_H\) and \(D_I\) are
evaluated.
\item \emph{Trace and centre.}  A bulk bound for \(\Dev_w(\Om')\) must be
converted into
\[
   \int_{\partial\Om'}
       \left|\nu-\frac{x-z}{|x-z|}\right|^2\,d\Hn
\]
for one centre \(z\), and that centre must also control the volume error in
the Faber--Krahn application.  This is an additional quantitative trace
statement, not a consequence of \eqref{eq:weighted-Reilly} alone.
\end{enumerate}

There is nevertheless one useful structural conclusion.  The level-set
moving-ball argument in the accompanying Faber--Krahn manuscript only needs
normal control after integration in the level parameter.  This is compatible
with the bulk nature of \(\Dev_w\).  If \(u\) is the ambient torsion function
and \(g=|\nabla u|\) on its regular level surfaces, the coarea formula shows
why the high-gradient issue survives integration:
\[
   \int dt\int_{\{u=t\}}g^3\,d\Hn
   =\int_{\Om}|\nabla u|^4\,dx.
\tag{5.23}\label{eq:coarea-fourth}
\]
Thus an integrated Reilly route still requires an \(L^4\)-type gradient
control, or an alternative estimate which avoids the cubic boundary moment.

\begin{problem}[A precise PDE-route target]\label{prob:pde-target}
Find hypotheses, valid for the torsion superlevels in the bounded
Faber--Krahn core, under which one can prove with computable constants:
\[
   \int_G\int_{\partial^*E_\rho}
      \left|\nu_\rho-\frac{x-z_\rho}{|x-z_\rho|}\right|^2
      d\Hn\,d\mu(\rho)
   \le C(n,R,\rho_*)\,\delta_T(\Om),
\tag{5.24}\label{eq:integrated-target}
\]
with the same centres \(z_\rho\) controlling the corresponding volume errors.
Theorem~\ref{thm:reilly-identities} is an exact identity available toward
this problem; (R1)--(R3) are the unresolved estimates.
\end{problem}

\section{Attempt IV: an intrinsic slicing and centre-motion route}
\label{sec:intrinsic}

The proof completed in Section~\ref{sec:selection} is effective, but its
organising principle is regularisation by penalised minimisers.  This
section records a different question: can one prove the strong
isoperimetric inequality by an argument whose variables directly exhibit
why a nearly optimal set is forced to be a single nearly spherical body?
The model is the Plan~2 torsion argument: a global deficit is decomposed
into levelwise geometric defects and a kinetic term paying for movement of
the comparison balls.

\subsection{The relevant precedent in quantitative isoperimetry}

The closest existing precedent is not the selection principle, but the
symmetrisation-and-induction proof of the sharp quantitative
isoperimetric inequality by Fusco--Maggi--Pratelli \cite{FMP}.  Their proof
passes through hyperplane sections, controls the loss incurred by
symmetrising those sections, and uses successive reflection symmetries to
remove the translation freedom of the comparison ball.  Thus it already
has the two features relevant here: an integrated family of
lower-dimensional deficits and a mechanism preventing independent centres
from drifting from slice to slice.

This is also the feature of the recent functional stability argument of
Figalli--van Hintum--Tiba \cite{FvHTPL}: in their
Pr\'ekopa--Leindler/Borell--Brascamp--Lieb proof, comparison translations
are first obtained on truncated ranges of level sets, and overlap between
successive ranges is then used to show that the translations can be chosen
coherently.  The analogy is structural, not yet a deduction: their
deficit is a functional Brunn--Minkowski deficit, whereas the quantity to
be controlled here is a boundary-normal defect.

\subsection{Why the fixed-centre form is the correct intrinsic quantity}

For \(z\in\R^n\) write
\[
   \mathcal B_z(E)
   :=\frac12\int_{\partial^*E}
        \left|\nu_E-\frac{x-z}{|x-z|}\right|^2\,d\Hn.
\tag{6.1}\label{eq:fixed-centre-B}
\]
Proposition~\ref{prop:radial-GG} gives, for \(|E|=|B_1|\),
\[
   \mathcal B_z(E)
   =\Per(E)-(n-1)\int_E\frac{dx}{|x-z|}
   =\Def(E)
     +(n-1)\left\{\int_{B_1(z)}\frac{dx}{|x-z|}
                       -\int_E\frac{dx}{|x-z|}\right\}.
\tag{6.2}\label{eq:B-def-plus-potential}
\]
The first term is the ordinary isoperimetric deficit.  The second is a
radial-displacement deficit about the \emph{same} centre.  Formula
\eqref{eq:B-def-plus-potential} explains why ordinary quantitative
isoperimetry alone is insufficient: small symmetric difference controls
from below, but not from above, the loss of a decreasing radial potential.
A small piece of mass transported far from the ball has radial-potential
loss of order its volume, even though its squared symmetric difference is
of smaller order.  A geometric FJ proof must show that the perimeter cost
of such displacement pays for this additional term.

\subsection{A rigorous slicing calculation}

Translate the proposed centre to \(0\), and write \(x=(s,y)\in
\R\times\R^{n-1}\).  Let \(E_s=\{y:(s,y)\in E\}\), \(v(s)=|E_s|\), and,
where \(v(s)>0\), let \(r(s)\ge0\) be defined by
\[
   v(s)=\omega_{n-1}r(s)^{n-1}.
\]
The transverse Schwarz symmetral of \(E\) is
\[
   E^\sharp:=\{(s,y):|y|<r(s)\}.
\tag{6.3}\label{eq:transverse-symmetral}
\]

\begin{lemma}[Potential monotonicity under transverse symmetrisation]
\label{lem:potential-transverse}
Let \(E\subset\R^n\) be measurable with finite measure and finite radial
potential about \(0\).  Then
\[
   \int_E\frac{dx}{|x|}
      \le \int_{E^\sharp}\frac{dx}{|x|}.
\tag{6.4}\label{eq:potential-transverse}
\]
If \(E\) also has finite perimeter, transverse Schwarz symmetrisation
does not increase perimeter.  Consequently,
\[
   \mathcal B_0(E)-\mathcal B_0(E^\sharp)
      =
      \bigl\{\Per(E)-\Per(E^\sharp)\bigr\}
      +(n-1)\left\{
          \int_{E^\sharp}\frac{dx}{|x|}
           -\int_E\frac{dx}{|x|}
       \right\}\ge0 .
\tag{6.5}\label{eq:B-symmetrisation-direction}
\]
\end{lemma}

\begin{proof}
For a.e.\ \(s\), the function
\(y\mapsto(s^2+|y|^2)^{-1/2}\) is radial and decreasing on
\(\R^{n-1}\).  Among sets of \((n-1)\)-dimensional measure \(v(s)\),
its integral is maximised by \(B_{r(s)}^{n-1}\).  Integrating the
resulting inequality in \(s\) proves \eqref{eq:potential-transverse}.
The perimeter assertion is the standard Schwarz symmetrisation
inequality for finite-perimeter sets, in the slicing form used in
\cite{FMP}.  Substitution in \eqref{eq:fixed-centre-B} proves
\eqref{eq:B-symmetrisation-direction}.
\end{proof}

The direction in \eqref{eq:B-symmetrisation-direction} is important.
Symmetrising produces a set whose oriented defect is smaller.  Therefore
one cannot merely symmetrise, prove FJ for \(E^\sharp\), and infer it for
\(E\).  One must estimate the oriented loss discarded by
symmetrisation.  This is exactly where a slice-motion term must enter.

For the axially symmetric output the oriented loss has a direct
one-dimensional expression.  Set
\(\sigma_{n-2}:=\mathcal H^{n-2}(S^{n-2})=(n-1)\omega_{n-1}\).

\begin{lemma}[Meridian action of an axially symmetric set]
\label{lem:meridian-action}
Let \(I=(a,b)\), let \(r\in W^{1,1}(I)\) be nonnegative with compact
support in \(I\), and suppose the set
\[
   F:=\{(s,y):s\in I,\ |y|<r(s)\}
\]
has finite perimeter and has no boundary portion other than the surface
generated by the graph \(r\).  Then
\[
\begin{aligned}
   \mathcal B_0(F)
   =\sigma_{n-2}\int_I r(s)^{n-2}
       \left[
          \sqrt{1+r'(s)^2}
          -\frac{r(s)-s r'(s)}
                 {\sqrt{s^2+r(s)^2}}
       \right]\,ds .
\end{aligned}
\tag{6.6}\label{eq:meridian-action}
\]
The integrand is nonnegative.  It vanishes a.e.\ exactly when the
outward normal of the generating curve agrees with its radial direction;
on each connected nontrivial component this means that the curve lies on
a circle centred at \(0\).
\end{lemma}

\begin{proof}
At a point \((s,r(s)\theta)\), \(\theta\in S^{n-2}\), the area element
and the two unit vectors are
\[
   d\Hn=\sigma_{n-2}r(s)^{n-2}\sqrt{1+r'(s)^2}\,ds,
\quad
   \nu_F=\frac{(-r'(s),\theta)}{\sqrt{1+r'(s)^2}},
\quad
   e_0=\frac{(s,r(s)\theta)}{\sqrt{s^2+r(s)^2}}.
\]
Since \(\tfrac12|\nu_F-e_0|^2=1-\nu_F\cdot e_0\), direct substitution
gives \eqref{eq:meridian-action}.  Its bracket is nonnegative by
Cauchy--Schwarz:
\[
   r-sr'\le\sqrt{s^2+r^2}\sqrt{1+(r')^2}.
\]
Equality means that \((-r',1)\) is a nonnegative multiple of
\((s,r)\); equivalently \(r'=-s/r\) wherever \(r>0\).  Hence
\((s^2+r(s)^2)'=0\) on each such component.
\end{proof}

Formula \eqref{eq:meridian-action} is the intrinsic analogue of the
Plan~2 action.  The profile \(r(s)\) is the moving radius of the
sections, and the integrand measures its failure to move according to
the profile of one fixed ball.  What has disappeared after transverse
symmetrisation is the motion of the section \emph{centres}; the
perimeter loss in \eqref{eq:B-symmetrisation-direction} must pay for
that lost motion.

\subsection{Attempt at the first lemma: centred reflections}

We now attempt the oriented symmetry reduction.  The basic operation in
the FMP reduction is to bisect a set by a hyperplane and reflect one of
the two halves.  For the oriented quantity there is a decisive
restriction: the reflecting hyperplane must contain the centre of the
radial field.  Under that restriction the operation has an exact
identity.

For \(q\in\R^n\) and a finite-measure set \(G\), write
\[
   I_q(G):=\int_G\frac{dx}{|x-q|}.
\]
The integral is finite: the kernel is integrable on \(B_1(q)\) and is
bounded by \(1\) outside \(B_1(q)\).

\begin{lemma}[One centred bisection]\label{lem:one-centred-bisection}
Let \(G\subset\R^n\) have finite perimeter and
\(|G|=\omega_n\).  Let \(q\in\R^n\), and let \(H\) be an affine
hyperplane containing \(q\), with complementary open half-spaces
\(H^+\) and \(H^-\).  Assume
\[
   |G\cap H^+|=|G\cap H^-|=\frac{\omega_n}{2}.
\tag{6.7a}\label{eq:centred-cut-assumptions}
\]
If \(R_H\) denotes reflection across \(H\), define
\[
   G^\pm:=(G\cap H^\pm)\cup R_H(G\cap H^\pm).
\]
Then \(|G^\pm|=\omega_n\), and
\[
\begin{aligned}
   \Def(G^+)+\Def(G^-)&=2\Def(G)-2\Per(G;H),\\
   \mathcal B_q(G^+)+\mathcal B_q(G^-)
      &=2\mathcal B_q(G)-2\Per(G;H).
\end{aligned}
\tag{6.7b}\label{eq:centred-cut-identities}
\]
Consequently there is a sign \(s\in\{+,-\}\) for which
\[
   \Def(G^s)\le2\Def(G),
   \qquad
   \mathcal B_q(G^s)\ge\mathcal B_q(G)-\Def(G).
\tag{6.7c}\label{eq:centred-cut-choice}
\]
\end{lemma}

\begin{proof}
Since \(H\) bisects \(G\) and reflection preserves Lebesgue measure,
\(|G^\pm|=2|G\cap H^\pm|=\omega_n\).

We first record the perimeter identity.  The reduced boundary of
\(G^+\) away from \(H\) consists of the reduced boundary of \(G\) in
\(H^+\) and its reflected copy.  Along \(H\) the two reflected traces
match, so there is no jump interface in \(G^+\).  The same observation
applies to \(G^-\).  Applying the BV restriction and gluing formula
on the two half-spaces gives
\[
   \Per(G^\pm)=2\Per(G;H^\pm),
   \qquad
   \Per(G^+)+\Per(G^-)=2\{\Per(G)-\Per(G;H)\}.
\tag{6.7d}\label{eq:perimeter-reflected-halves}
\]
All four sets in this equality have volume \(\omega_n\); subtracting
\(2\Per(B_1)\) proves the first identity in
\eqref{eq:centred-cut-identities}.

Because \(q\in H\), reflection fixes \(q\) and therefore
\(|R_Hx-q|=|x-q|\).  It follows that
\[
   I_q(G^\pm)=2I_q(G\cap H^\pm),
   \qquad I_q(G^+)+I_q(G^-)=2I_q(G).
\tag{6.7e}\label{eq:potential-reflected-halves}
\]
Insert \eqref{eq:perimeter-reflected-halves} and
\eqref{eq:potential-reflected-halves} in
\(\mathcal B_q(A)=\Per(A)-(n-1)I_q(A)\), which is
\eqref{eq:fixed-centre-B}, to obtain the second identity in
\eqref{eq:centred-cut-identities}.

By the isoperimetric inequality both deficits in the first line of
\eqref{eq:centred-cut-identities} are nonnegative.  Consequently
\[
   \Per(G;H)\le\Def(G),
\]
and each reflected deficit is at most \(2\Def(G)\).  By the second
line at least one of
\(\mathcal B_q(G^+)\), \(\mathcal B_q(G^-)\) is at least
\(\mathcal B_q(G)-\Per(G;H)\), hence at least
\(\mathcal B_q(G)-\Def(G)\).  Choose that sign.
\end{proof}

\begin{remark}[What a centred facet costs]\label{rem:facet-loss}
The terms involving \(\Per(G;H)\) in
\eqref{eq:centred-cut-identities} have a direct meaning.  The
reflected constructions erase every reduced-boundary
portion lying in \(H\).  On such a portion \(e_q(x)\) is tangent to
\(H\), while \(\nu_G(x)\) is normal to \(H\), so
\(\tfrac12|\nu_G-e_q|^2=1\): its oriented cost is exactly its
perimeter.  Crucially, this erased cost is not an obstruction at the
level of one bisection: nonnegativity of the two reflected deficits
automatically pays it through \(\Per(G;H)\le\Def(G)\).
\end{remark}

The next theorem iterates Lemma~\ref{lem:one-centred-bisection}.
It proves every centred reflection which can be forced by an
intermediate-value argument.

\begin{theorem}[Centred transverse symmetry reduction]
\label{thm:transverse-oriented-reduction}
Let \(E\subset\R^n\), \(n\ge2\), be a finite-perimeter set with
\(|E|=\omega_n\), and let \(q\in\R^n\).
There exist \(n-1\) mutually orthogonal hyperplanes
\(H_1,\ldots,H_{n-1}\), all containing \(q\), and a set
\(F\subset\R^n\), symmetric with respect to every \(H_i\), such that
\[
   |F|=\omega_n,\qquad
   \Def(F)\le2^{\,n-1}\Def(E),\qquad
   \mathcal B_q(E)
      \le \mathcal B_q(F)+(2^{\,n-1}-1)\Def(E).
\tag{6.7g}\label{eq:transverse-reduction-output}
\]
In particular, if \(q\) is a minimising centre for \(\beta(E)\), then,
after translating \(q\) to \(0\) and rotating the line
\(\bigcap_{i=1}^{n-1}H_i\) to the \(x_1\)-axis,
\[
   \beta(E)^2
      \le\mathcal B_0(F)+(2^{\,n-1}-1)\Def(E),
   \qquad
   \Def(F)\le2^{\,n-1}\Def(E),
\tag{6.7h}\label{eq:transverse-reduction-beta}
\]
and \(F\) is symmetric in all transverse coordinate variables.
\end{theorem}

\begin{proof}
Set \(F_0=E\).  We inductively construct \(F_j\) and mutually
orthogonal unit normals \(\nu_1,\ldots,\nu_j\), for
\(0\le j\le n-1\), such that \(F_j\) is symmetric with respect to
\[
   H_i:=q+\nu_i^\perp,\qquad 1\le i\le j,
\]
and
\[
   \Def(F_j)\le2^j\Def(E),
   \qquad
   \mathcal B_q(F_j)
      \ge\mathcal B_q(E)-(2^j-1)\Def(E).
\tag{6.7i}\label{eq:induction-reflections}
\]

Assume \(j\le n-2\) and \(F_j\) has been constructed.  Put
\[
   V_j:=\{\nu_1,\ldots,\nu_j\}^{\perp};
\]
then \(\dim V_j=n-j\ge2\).  For \(\nu\in S(V_j)\) define
\[
   \Phi_j(\nu)
   :=|F_j\cap\{(x-q)\cdot\nu>0\}|
      -|F_j\cap\{(x-q)\cdot\nu<0\}|.
\]
Since hyperplanes have zero \(n\)-dimensional Lebesgue measure,
dominated convergence shows that \(\Phi_j\) is continuous.  Moreover
\(\Phi_j(-\nu)=-\Phi_j(\nu)\).  The sphere \(S(V_j)\) is connected, so
there is \(\nu_{j+1}\in S(V_j)\) with
\(\Phi_j(\nu_{j+1})=0\).  Thus
\(H_{j+1}=q+\nu_{j+1}^{\perp}\) bisects \(F_j\).

Apply Lemma~\ref{lem:one-centred-bisection} and select the sign in
\eqref{eq:centred-cut-choice}; call the selected reflected set
\(F_{j+1}\).  Since \(\nu_{j+1}\perp\nu_i\), reflection across
\(H_{j+1}\) commutes with reflection across each \(H_i\), and each of
the two half-spaces bounded by \(H_{j+1}\) is invariant under the
previous reflections.  Thus \(F_{j+1}\) retains all previous
symmetries and gains the new one.  Estimate
\eqref{eq:induction-reflections} follows because
\[
\begin{aligned}
   \mathcal B_q(F_{j+1})
   &\ge\mathcal B_q(F_j)-\Def(F_j)\\
   &\ge\mathcal B_q(E)
       -\{(2^j-1)+2^j\}\Def(E)
    =\mathcal B_q(E)-(2^{j+1}-1)\Def(E).
\end{aligned}
\]
At \(j=n-1\) this proves
\eqref{eq:transverse-reduction-output}; invariance under rigid motions
and \(\mathcal B_q(E)=\beta(E)^2\) give
\eqref{eq:transverse-reduction-beta}.
\end{proof}

\begin{remark}[Why the iteration stops at \(n-1\)]
\label{rem:last-centred-reflection}
After \(j\) reflections the new normal is searched for in
\(V_j=\{\nu_1,\ldots,\nu_j\}^{\perp}\).  Up to \(j=n-2\), the sphere
\(S(V_j)\) is connected and oddness of \(\Phi_j\) produces a
centred bisection.  At \(j=n-1\), \(V_j\) is one-dimensional:
\(S(V_j)=\{\nu,-\nu\}\), and oddness gives no zero.  The last
orthogonal plane through \(q\) need not bisect \(F_{n-1}\).  Moving
that plane until it becomes a median is the step used in the FMP
volume-asymmetry reduction, but it no longer fixes \(q\), so
\(|R_Hx-q|=|x-q|\) and \eqref{eq:potential-reflected-halves} fail.
This is precisely the axial centre-motion term still missing from the
oriented argument.
\end{remark}

\begin{proposition}[The exact last-axis target]\label{prop:last-axis-target}
Let \(q\) be a minimising centre for \(\beta(E)\), and let \(G\) be
the set produced by Theorem~\ref{thm:transverse-oriented-reduction},
so that the intersection of its \(n-1\) symmetry hyperplanes is a
line \(L\) through \(q\).  Let \(c\in L\) be chosen so that the
hyperplane \(H_c\) through \(c\), perpendicular to \(L\), bisects
\(G\).  Suppose that, for some dimensional number \(K_n\),
\[
   (n-1)\bigl\{I_c(G)-I_q(G)\bigr\}
      \le K_n\Def(E).
\tag{6.7j}\label{eq:last-axis-potential-shift}
\]
Then there is an \(n\)-symmetric set \(F\), centred at \(c\), with
\[
   \Def(F)\le2^n\Def(E),
   \qquad
   \beta(E)^2
      \le \mathcal B_c(F)
           +\{2^n-1+K_n\}\Def(E).
\tag{6.7k}\label{eq:last-axis-output}
\]
\end{proposition}

\begin{proof}
The fixed-centre identity gives
\[
   \mathcal B_q(G)-\mathcal B_c(G)
      =(n-1)\{I_c(G)-I_q(G)\}
      \le K_n\Def(E).
\]
Apply Lemma~\ref{lem:one-centred-bisection} to \(G\), now with centre
\(c\) and cutting plane \(H_c\).  Choose one reflected half, denoted
by \(F\), for which
\[
   \Def(F)\le2\Def(G),
   \qquad
   \mathcal B_c(G)\le\mathcal B_c(F)+\Def(G).
\]
Every earlier symmetry hyperplane contains the axis \(L\), hence
contains \(c\); it is orthogonal to \(H_c\).  Its reflection commutes
with reflection across \(H_c\), so \(F\) has all \(n\) symmetries and
their intersection is \(c\).  Finally,
\[
\begin{aligned}
   \beta(E)^2
   &\le \mathcal B_q(G)+(2^{n-1}-1)\Def(E)\\
   &\le \mathcal B_c(F)+\Def(G)
        +\{K_n+2^{n-1}-1\}\Def(E)\\
   &\le \mathcal B_c(F)+\{K_n+2^n-1\}\Def(E),
\end{aligned}
\]
while \(\Def(F)\le2\Def(G)\le2^n\Def(E)\).
\end{proof}

\begin{proposition}[What a full orthogonal equipartition would give]
\label{prop:orthant-oriented-reduction}
Let \(E\) have finite perimeter and \(|E|=\omega_n\), and let
\(q\in\R^n\).  Suppose that there are \(n\) mutually orthogonal
hyperplanes through \(q\) whose \(2^n\) open orthants \(Q_\sigma\)
satisfy
\[
   |E\cap Q_\sigma|=2^{-n}\omega_n
   \quad\hbox{for every \(\sigma\in\{+,-\}^n\)}.
\tag{6.7l}\label{eq:orthant-equipartition}
\]
Then there is an \(n\)-symmetric set \(F\), with symmetry centre \(q\),
such that
\[
   |F|=\omega_n,\qquad
   \Def(F)\le2^n\Def(E),\qquad
   \mathcal B_q(E)\le\mathcal B_q(F)+\Def(E).
\tag{6.7m}\label{eq:full-orthant-output}
\]
\end{proposition}

\begin{proof}
For each orthant \(Q_\sigma\), reflect \(E\cap Q_\sigma\) through the
group generated by the \(n\) hyperplane reflections, and denote the
resulting \(n\)-symmetric set by \(F_\sigma\).  By
\eqref{eq:orthant-equipartition}, \(|F_\sigma|=\omega_n\).  Put
\(\Sigma:=\bigcup_{i=1}^nH_i\).  As in
Lemma~\ref{lem:one-centred-bisection}, boundary on \(\Sigma\) is
erased by reflection and the radial kernel about \(q\) is invariant
under all the reflections.  Summing over the \(2^n\) orthants gives
\[
   \sum_\sigma\Def(F_\sigma)
      =2^n\{\Def(E)-\Per(E;\Sigma)\},
   \qquad
   \sum_\sigma\mathcal B_q(F_\sigma)
      =2^n\{\mathcal B_q(E)-\Per(E;\Sigma)\}.
\]
Every deficit is nonnegative, so
\(\Per(E;\Sigma)\le\Def(E)\), and every individual deficit is at most
\(2^n\Def(E)\).  Choose an orthant for which
\(\mathcal B_q(F_\sigma)\ge
\mathcal B_q(E)-\Per(E;\Sigma)\); it satisfies the claimed estimates.
\end{proof}

\begin{lemma}[Orthogonal quadrisection in the plane]
\label{lem:orthogonal-quadrisection}
Let \(E\subset\R^2\) be measurable with \(0<|E|<\infty\).  There are
two perpendicular lines \(H_1,H_2\) whose four open quadrants each
contain measure \(|E|/4\) of \(E\).
\end{lemma}

\begin{proof}
We first suppose that the measure \(\mu=\mathbf1_E\,dx\) is replaced
by a finite measure having a strictly positive continuous density and
finite total mass \(m\).  For every \(\theta\in\R\), put
\[
   u_\theta=(\cos\theta,\sin\theta),
   \qquad v_\theta=(-\sin\theta,\cos\theta).
\]
There is a unique number \(a(\theta)\) such that the oriented line
\[
   H_\theta:=\{x:x\cdot u_\theta=a(\theta)\}
\]
bisects \(\mu\); positivity and continuity of the density imply that
\(a(\theta)\) depends continuously on \(\theta\).  There is likewise
a unique \(b(\theta)\), continuous in \(\theta\), such that
\[
   K_\theta:=\{x:x\cdot v_\theta=b(\theta)\}
\]
bisects \(\mu\).  Let
\[
   Q_\theta^{++}
   :=\{x:x\cdot u_\theta>a(\theta),\
          x\cdot v_\theta>b(\theta)\},
   \qquad
   f(\theta):=\mu(Q_\theta^{++})-\frac m4.
 \]
The function \(f\) is continuous.  At angle \(\theta+\pi/2\), the
two bisecting lines are the same lines with the roles of their
positive sides changed:
\[
   Q_{\theta+\pi/2}^{++}
   =\{x:x\cdot v_\theta>b(\theta),\
          x\cdot u_\theta<a(\theta)\}.
 \]
Since \(K_\theta\) bisects \(\mu\), this gives
\[
   f(\theta+\pi/2)
   =\left\{\frac m2-\mu(Q_\theta^{++})\right\}-\frac m4
   =-f(\theta).
 \]
The intermediate value theorem provides \(\theta_0\) with
\(f(\theta_0)=0\).  Both lines bisect the measure and one quadrant
has mass \(m/4\), so each of the remaining three quadrants also has
mass \(m/4\).

For the original measure, take smooth strictly positive densities
\(\mu_k\) converging narrowly to \(\mathbf1_E\,dx\), with total mass
\(|E|\), for example by convolving with a Gaussian of variance
\(1/k\).  Choose perpendicular quadrisecting lines for \(\mu_k\).
Their normals have a convergent subsequence.  Their offsets remain
bounded: otherwise one of the corresponding half-spaces would contain
mass tending to \(0\) or to \(|E|\), although it contains
\(|E|/2\) for every \(k\).  Pass to a subsequence so both lines
converge.  A line has two-dimensional Lebesgue measure zero, hence
the four limiting open quadrants are continuity sets for
\(\mathbf1_E\,dx\).  Narrow convergence therefore passes the
quadrant masses to the limit and gives \(|E|/4\) in each quadrant.
\end{proof}

\begin{corollary}[Full oriented reduction in dimension two]
\label{cor:oriented-reduction-plane}
Let \(E\subset\R^2\) have finite perimeter and \(|E|=\omega_2\).
There is a set \(F\), symmetric with respect to two perpendicular
lines meeting at \(q\), such that, after translating \(q\) to \(0\),
\[
   |F|=\omega_2,\qquad
   \Def(F)\le4\Def(E),\qquad
   \beta(E)^2\le\mathcal B_0(F)+\Def(E).
\tag{6.7n}\label{eq:planar-oriented-reduction}
\]
\end{corollary}

\begin{proof}
Apply Lemma~\ref{lem:orthogonal-quadrisection} and let \(q\) be the
intersection of the two quadrisecting lines.  Proposition
~\ref{prop:orthant-oriented-reduction}, with \(n=2\), gives
\(\mathcal B_q(E)\le\mathcal B_q(F)+\Def(E)\) and the asserted
deficit estimate.  Since
\(\beta(E)^2=\inf_z\mathcal B_z(E)\le\mathcal B_q(E)\), translation
completes the proof.
\end{proof}

Proposition~\ref{prop:orthant-oriented-reduction} completely proves
the desired oriented reduction if an orthogonal equipartition is
available.  Corollary~\ref{cor:oriented-reduction-plane}
therefore closes the first missing lemma in dimension \(2\).  This is
not a dimension-free device:
Maldonado and Rold\'an-Pensado \cite{MaldonadoRoldan} construct, for
every \(n\ge3\), measures in \(\R^n\) admitting no equipartition into
\(2^n\) equal parts by \(n\) mutually orthogonal hyperplanes.  Thus
the higher-dimensional proof cannot obtain its final centred symmetry
merely by reflecting one equal-volume orthant.

\subsection{The three geometric lemmas that would close this route}

The FMP proof suggests the following route, which uses no penalised
minimisation and no regularity theorem.

\begin{problem}[Full oriented symmetry reduction]\label{prob:oriented-reduction}
Remove the remaining restrictions in
Theorem~\ref{thm:transverse-oriented-reduction}: the resulting set
should acquire the last axial symmetry in a manner compatible with
the same-centre quantity needed
in Problem~\ref{prob:meridian-closure}.  More precisely, construct from arbitrary \(E\) an
\(n\)-symmetric set \(F\) whose symmetry centre is \(0\), with
\[
   \Def(F)\le C_n\Def(E),
   \qquad
   \beta(E)^2\le C_n\bigl\{\Def(E)+\mathcal B_0(F)\bigr\}.
\tag{6.7}\label{eq:oriented-reduction-target}
\]
The centred reflection theorem proves the transverse \(n-1\)
symmetries without losing oriented defect.  The final FMP reflection
uses an axial median hyperplane, which need not contain the tracked
radial centre; by Proposition~\ref{prop:last-axis-target}, it remains
to prove the potential-shift bound
\eqref{eq:last-axis-potential-shift}.  This is the unresolved
centre-motion step.
\end{problem}

\begin{problem}[Transverse slice-motion estimate]\label{prob:slice-motion}
For an \(n\)-symmetric, small-deficit, bounded set \(F\), prove
\[
   \mathcal B_0(F)-\mathcal B_0(F^\sharp)
      \le C_n\bigl\{\Per(F)-\Per(F^\sharp)\bigr\}.
\tag{6.8}\label{eq:slice-motion-target}
\]
In a nonsymmetrised formulation, the left side should be represented by
the variation of a section-centre field \(a(s)\); the estimate would
then be the isoperimetric analogue of controlling the centre velocity
in Plan~2.
\end{problem}

\begin{problem}[Centred meridian closure]\label{prob:meridian-closure}
For axially symmetric, volume-normalised \(F=F^\sharp\), assume in
addition that \(0\) is a radial-potential maximiser for \(F\), namely
\[
   I_0(F)=\sup_{z\in\R^n}I_z(F).
\]
Prove directly from \eqref{eq:meridian-action} that
\[
   \mathcal B_0(F)
      \le C_n\{\Per(F)-\Per(B_1)\}.
\tag{6.9}\label{eq:meridian-target}
\]
This is a one-dimensional weighted inequality for the generating
profile \(r\).  It is the first part of the intrinsic programme to
attempt once the centre is inherited correctly.  The centring
assumption cannot be omitted: \(F=B_1(ae_1)\), \(a\ne0\), is axially
symmetric and has zero perimeter deficit, whereas
\(\mathcal B_0(F)>0\).
\end{problem}

If \eqref{eq:oriented-reduction-target}--\eqref{eq:meridian-target}
hold, their composition proves Fusco--Julin by a geometric
symmetrisation argument.  More importantly, the proof would expose
the mechanism: transverse perimeter pays for movement of section
centres, while meridian perimeter pays for failure of the section
radii to follow one centred sphere.

\subsection{Sobolev and Minkowski formulations: useful diagnostics, not yet proofs}

The \(BV\)-Sobolev equivalence does contain an exact integrated-level
identity.  For \(q=n/(n-1)\) and a nonnegative compactly supported
\(f\in BV(\R^n)\), write \(E_t=\{f>t\}\).  Coarea and Minkowski's
integral inequality give
\[
\begin{aligned}
   |Df|(\R^n)-n\omega_n^{1/n}\|f\|_{L^q}
   &=
   \int_0^\infty
     \left\{\Per(E_t)-n\omega_n^{1/n}|E_t|^{(n-1)/n}\right\}\,dt \\
   &\quad+
   n\omega_n^{1/n}
   \left\{\int_0^\infty |E_t|^{(n-1)/n}\,dt-\|f\|_{L^q}\right\},
\end{aligned}
\tag{6.10}\label{eq:BV-Sobolev-splitting}
\]
and both terms on the right are nonnegative.  Thus Sobolev does
produce integrated isoperimetric deficits.  It does not, by itself,
solve the present problem.  If \(f=\mathbf1_E\), there is only one
nontrivial level and no centre-motion information.  If \(f\) is a
smoothing or a potential associated to \(E\), one still needs a
non-circular estimate bounding its Sobolev deficit by \(\Def(E)\) and
a trace statement recovering \(\mathcal B_z(E)\).  Applying
quantitative isoperimetry or Fusco--Julin separately to the levels
would assume the desired input rather than derive it.

Likewise, outer parallel sets connect isoperimetry to Brunn--Minkowski:
for a sufficiently regular set,
\[
   \left.\frac{d}{dt}\right|_{t=0^+}|E+tB_1|=\Per(E).
\tag{6.11}\label{eq:parallel-first-var}
\]
The stability results of Figalli--van Hintum--Tiba
\cite{FvHTBM,FvHTPL} make this connection compelling at the level of
volume and coherent translations.  However, their conclusions do not
presently bound the boundary direction term
\(\mathcal B_z(E)\).  An isoperimetric application would require a new
oriented first-variation version of Brunn--Minkowski stability.

\begin{remark}[Current assessment]\label{rem:intrinsic-assessment}
The slicing route is the most faithful analogue of Plan~2 presently
visible in the isoperimetric literature.  It starts from a proof method
already designed to prevent translation drift, and
\eqref{eq:meridian-action} shows exactly where the FJ normal
oscillation sits.  The Sobolev and Brunn--Minkowski viewpoints help
organise possible identities, but neither removes the centre-motion
lemma without an additional, presently unproved, oriented estimate.
\end{remark}

\section{Consequences for the Faber--Krahn manuscript}\label{sec:application}

The level-set moving-ball proof needs a same-centre estimate: the centre used
to compare \(E_\rho\) with \(B_\rho\) in measure must also control the radial
normal error.  Corollary~\ref{cor:scaled-graph} has exactly that form once a
balanced graph representation is available.

At present, the following statements are justified.
\begin{enumerate}[label=\textnormal{(\arabic*)},leftmargin=2.6em]
\item The black-box Fusco--Julin import gives the desired normal control on
all finite-perimeter levels, but its dimensional constant is not made
computable by the published proof.
\item Theorem~\ref{thm:graph-main} and Corollary~\ref{cor:scaled-graph} give a
fully explicit replacement on balanced nearly-spherical levels.
\item The potential-deficit selection
Theorem~\ref{thm:FJ-selection} preserves \(\Def/\mathcal C\), but its
near-minimality error has perimeter order.  The regularising replacement is
Theorem~\ref{thm:true-beta-selection}: together with
Theorem~\ref{thm:additive-almost-min}, it preserves
\(\Def/\beta^2\) and produces an additive perimeter almost-minimiser with
displayed constants.  Proposition~\ref{prop:flat-number-computable}
certifies computability of the graph-entry threshold, and
Proposition~\ref{prop:true-beta-bounded-close} closes the resulting
strong inequality in the bounded class with computable constants.
\item Lemma~\ref{lem:computable-confinement} gives the computable
confinement reduction for arbitrary finite-perimeter sets.  Consequently
Theorem~\ref{thm:global-computable-FJ} proves its normal-oscillation form,
and Corollary~\ref{cor:original-FJ-computable} proves the original
same-centre combined formulation, with computable dimensional constants.
\item The weighted Reilly calculations prove exact identities tailored to
torsion levels, but do not yet provide the missing normal estimate in the
generality of arbitrary bounded open sets.
\item Section~\ref{sec:intrinsic} formulates a separate intrinsic programme.
It proves the fixed-centre symmetrisation and meridian-action identities
and Theorem~\ref{thm:transverse-oriented-reduction}, which performs the
oriented reflection reduction through \(n-1\) transverse symmetries.
Corollary~\ref{cor:oriented-reduction-plane} supplies the full first
reduction in dimension \(2\).  In dimensions \(n\ge3\), the last
centred symmetry, transverse slice-motion control, and centred
one-dimensional meridian closure remain open in that route.
\end{enumerate}

Accordingly, the regularised-selection branch provides a truly quantitative
proof of the Fusco--Julin inequality: all constants entering
\eqref{eq:global-computable-FJ} and \eqref{eq:original-FJ-computable} are
computable, although no numerical value is claimed.  The integrated PDE
target in Problem~\ref{prob:pde-target} remains a separate possible route
for the Faber--Krahn application.

\section{Checklist for future attempts}\label{sec:checklist}

Any claimed effective proof of the general FJ inequality should state, and
permit a reader to verify, each of the following items.
\begin{enumerate}[label=\textnormal{(\arabic*)},leftmargin=2.6em]
\item The exact convention for \(\beta(E)^2\), including the factor \(1/2\).
\item The scaling of both \(\beta(E)^2\) and \(\Def(E)\).
\item The centre used for every same-centre conclusion.
\item A displayed small-deficit threshold and the trivial large-deficit
branch.
\item If nearly-spherical analysis is used, the handling of the degree-zero
and degree-one spherical harmonic modes.
\item If a regularisation theorem is used, computable almost-minimality and
regularity parameters.
\item If a torsion identity is used, the sign convention
\(-\Delta v=1\), \(v>0\), \(v=0\) on the boundary, and the estimates needed
to pass from a bulk Hessian defect to a boundary normal defect.
\end{enumerate}

The calculations in this article settle items (1)--(5) in the
nearly-spherical regime, and Theorem~\ref{thm:additive-almost-min}
and Proposition~\ref{prop:flat-number-computable} settle item (6) for
the bounded regularised-selection route; Lemma~\ref{lem:computable-confinement}
globalises that route.  The estimates in item (7) remain open only for the
separate PDE route.

\begin{thebibliography}{99}

\bibitem{FJ}
N.~Fusco and V.~Julin,
\emph{A strong form of the quantitative isoperimetric inequality},
Calc. Var. Partial Differential Equations \textbf{50} (2014), 925--937.
Preprint: \url{https://arxiv.org/abs/1111.4866}.

\bibitem{FMP}
N.~Fusco, F.~Maggi, and A.~Pratelli,
\emph{The sharp quantitative isoperimetric inequality},
Ann. of Math. (2) \textbf{168} (2008), 941--980.

\bibitem{FigalliMaggiPratelli}
A.~Figalli, F.~Maggi, and A.~Pratelli,
\emph{A mass transportation approach to quantitative isoperimetric
inequalities},
Invent. Math. \textbf{182} (2010), 167--211.

\bibitem{FvHTBM}
A.~Figalli, P.~van Hintum, and M.~Tiba,
\emph{Sharp quantitative stability of the Brunn--Minkowski inequality},
preprint, 2023.
\url{https://arxiv.org/abs/2310.20643}.

\bibitem{FvHTPL}
A.~Figalli, P.~van Hintum, and M.~Tiba,
\emph{Sharp quantitative stability for the Pr\'ekopa--Leindler and
Borell--Brascamp--Lieb inequalities},
preprint, 2025.
\url{https://arxiv.org/abs/2501.04656}.

\bibitem{MaldonadoRoldan}
G.~L.~Maldonado and E.~Rold\'an-Pensado,
\emph{On the orthogonal Gr\"unbaum partition problem in dimension three},
preprint, 2024.
\url{https://arxiv.org/abs/2404.01504}.

\bibitem{Fuglede}
B.~Fuglede,
\emph{Stability in the isoperimetric problem for convex or nearly spherical
domains in \(\R^n\)},
Trans. Amer. Math. Soc. \textbf{314} (1989), 619--638.

\bibitem{FuscoSurvey}
N.~Fusco,
\emph{The quantitative isoperimetric inequality and related topics},
Bull. Math. Sci. \textbf{5} (2015), 517--607.
\url{https://doi.org/10.1007/s13373-015-0074-x}.

\bibitem{Maggi}
F.~Maggi,
\emph{Sets of Finite Perimeter and Geometric Variational Problems},
Cambridge Studies in Advanced Mathematics, vol.~135,
Cambridge University Press, 2012.

\end{thebibliography}

\end{document}
